\chapter{流形及其张量场}\label{appx:manifold}

最初仅用微积分和坐标系研究“背景”的\textbf{局部}(local)性质，建立了\textbf{微分几何}(differential geometry)。随着集合论和拓扑学发展，我们发现可将“背景”准确定义为拓扑空间，且拓扑学提供了分析\textbf{整体}(global)性质的工具，便有了\textbf{现代微分几何}。

\section{微分流形}
物理背景作为拓扑空间，是无特征且无定形的，其中的点和曲线是不能先验识别或定义的，除非利用坐标系。有用的背景从来不可能真正脱离坐标。
坐标方法就是一种给背景点连续指定 $\F^n$ 数据的映射。
例如，平面上附着不同的坐标空间，一种可存放正交坐标，一种可存放极坐标，从而可赋予某点 $p$ 正交坐标 $(0,1)$，也可赋予极坐标 $(1,\pi/2)$，如图 \ref{fig:coor-chan}。拓扑学的同胚关系可描述坐标的指定，因为同胚具有双向连续性，且两个点不会被分配到同一坐标，从而避免歧义。
当然，同胚使得许多坐标系能力有限，不能揽下全局。比如，极坐标系不能完整覆盖 $\R^2$，坐标取值为 $V=(0,+\infty)\times(0,2\pi)$。因此常考察开子集。为圆满完成指派任务，整个背景须由若干坐标系拼接，且互相分担对方力所不及之地。本节主要考虑 $\R^n$。


\begin{definition}
    置 $(M,\topo T)$。存在 $n\in\N$，使 $M$ 任意点都存在邻域同胚于 $(\R^n,\topo T_u)$ 某开集，则称其\textbf{局部欧}。
    $\psi$ 称为\textbf{坐标映射}，$\psi(p)$ 称为 $p$ 的\textbf{坐标}。$O$ 称为\textbf{坐标域}(patch)，$(O,\psi)$ 称为\textbf{坐标系}或\textbf{图}(chart)。相应的坐标系族 $\mathcal A=\{(O,\psi):O\in\topo P\}$ 称为\textbf{图册}(atlas)。$(M,\mathcal A)$ 称为\textbf{局部欧氏空间}。维数为 $\dim M:=n$。当 $\mathcal A$ 自明时简记 $M$。
\end{definition}

\begin{remark}
    同胚连续性由全集拓扑或诱导拓扑衡量。原则上可修改定义，使其允许各坐标系不同维，这样空间无唯一维数，但实用价值不大。
\end{remark}

\begin{theorem}
    局部欧的等价说法是任意点存在邻域同胚于开球或 $\R^n$。
\end{theorem}
\begin{proof}
    因为 $\psi(x)$ 的邻域内一定存在包含该点的开球。再注意开球与 $\R^n$ 同胚。
\end{proof}

\begin{figure}[ht]\centering
    \includegraphics[width=.8\textwidth]{fig/appx/Europe-Blank.png}
    \caption{地图、地图册}
\end{figure}

同胚赋予的第 $i$ 个坐标构成映射 $x^i:O\to\R$，坐标系又常记 $\{x\}$。
就某个坐标系而言，坐标域的子集
    $\{p\in O: x^i(p)$ 是常函数，$i=2,\cdots,n\}$
可看成以 $x^1$ 为参数的路径，称其是一条 $x^1$ \textbf{坐标线}(coordinate line)。其它情况以此类推。全体坐标线构成\textbf{坐标网格}。
坐标域是某点邻域，故坐标系又称局部坐标系，同胚又称局部同胚。仅用一个坐标域就能覆盖则称\textbf{平凡的}(trivial)，此时可称坐标系是整体坐标系。典型的例子是 $\R^n$ 自身。设 $a=(x^1_a,\cdots,x^n_a)\in\R^n$，恒等映射给出\textbf{自然坐标} $\psi(a)=(x^1(a),\cdots,x^n(a))=a$。

\begin{definition}
    仿紧 $T_2$ 局部欧氏空间称为\textbf{拓扑流形}，简记 $M$。将 $\R^n$ 换为 $\C^n$ 就是\textbf{复流形}。
\end{definition}

\begin{remark}
    仿紧 $T_2$ 是为了应用单位分解。
    也可用更严苛的 $C_2,T_2$ 条件。
    下一节将说明 $C_2\land T_2\land\text{局部欧}\im \text{仿紧}\land T_2\land\text{局部欧} \im \text{可度量}$。
\end{remark}

\begin{figure}[ht]
    \centering
\includegraphics[width=.4\textwidth]{fig/appx/coor-chan.pdf}
    \caption{坐标变换}\label{fig:coor-chan}
\end{figure}

由于坐标域都是开的，故必存在相交之地。此处便涉及坐标系的转化。

\begin{definition}
    取 $(O_\alpha,\psi_\alpha),(O_\beta,\psi_\beta)\in \mathcal A$。若 $O_\alpha\cap O_\beta=\mt$；或 $O_\alpha\cap O_\beta\ne\mt$ 时\textbf{坐标变换}或\textbf{转移函数} $\psi_\alpha\circ\psi_\beta^{-1},\psi_\beta\circ\psi_\alpha^{-1}$ 的 $r$ 阶导函数连续（记 $C^r$），则称两坐标系 $C^r$\textbf{-相容}。
    所含坐标系皆 $C^r$-相容的 $\mathcal A$ 称为 $C^r$\textbf{-图册}。$C^\infty$-相容简称\textbf{相容}。
\end{definition}

\begin{remark}
    坐标变换可显式地表为 $n$ 个 $n$ 元函数，$C^r$ 性由 $\R^n,\C^n$ 分析学给出。$C^0$ 即连续。光滑即任意阶导函数连续，记 $C^\infty$。解析是指任意点可展开为收敛于自身的幂级数，从而一定光滑。由于光滑仍可做到在某点无穷阶平坦，如 $y=\e^{-1/x}$ 于 $x=0$（左侧 $x\leqslant 0$ 连接 $y=0$），因而从该点展开收敛不到自身，即光滑不一定解析。由此可构造实用的\textbf{隆起函数}(bump function)
    \eq{
        \Psi(x)=\begin{cases}
\exp(\frac{1}{x^2-1}), & |x|<1,\\
0, & |x|\geqslant 1.
\end{cases}
    }
    考虑复变函数 $f(z)=u(z)+\i v(z)$ 对定义域内 $z=x+\i y$ 可导，即 $f'(z):=\lim_{\Delta z\to 0}(f(z+\Delta z)-f(z))/\Delta z$ 需要对任意 $\Delta z=\Delta x+\i\Delta y$ 收敛。只需分别取 $\Delta x,\Delta y=0$ 就有 $\pdv{u}{x}+\i\pdv{v}{x}=\frac{1}{\i}\left(\pdv{u}{y}+\i\pdv{v}{y}\right)$，即 Cauchy-Riemann 条件
    \eq{
        \pdv{u}{x} = \pdv{v}{y},\quad \pdv{v}{x} = -\pdv{u}{y}.
    }
    因此虽然 $\C,\R^2$ 同胚，但 $\C$ 多出来乘法、线性等结构使复可导要比实可导多出上述条件，从而可证复变函数只需可导（称为全纯）就等价于解析。
\end{remark}

\begin{example}
    拓扑流形相当于仅要求 $C^0$-相容。
\end{example}

假设有两个图册 $\mathcal A,\mathcal B$。它们可能并不相容，即其间存在相交但不相容的坐标系。如果它们相容，则可合并为一个更大的图册。当然希望将所有相容的坐标系都囊括进来，尽可能构成足够大的图册。

\begin{definition}
    一个不真含于任何 $C^r$-图册的 $C^r$-图册称为\textbf{最大的}(maximal)。若局部欧氏空间的图册是 $C^r$-最大图册，则称 \textbf{$C^r$-流形}。最大图册又称\textbf{微分结构}。$C^\infty$-流形称为\textbf{微分流形}或\textbf{光滑流形}。
\end{definition}

\begin{example}
    \textbf{通常流形}就是指 $(\R^n,\topo T_u)$ 赋以包含自然坐标系的最大光滑图册。
\end{example}

综上，将今后常用的 $C^\infty$-拓扑流形之定义总结为：
\begin{tcolorbox}
\begin{definition}
    置 $T_2$ 仿紧空间 $(M,\topo T)$，若 $\exists\topo P\subset\topo T$ 使
    \begin{itemize}
        \item $\topo P$ 是开覆盖：$M = \bigcup_{O\in\topo P}O$；
        \item $\forall O\in\topo P$，$\exists$ 同胚 $\psi:O\to V\in\topo T_u$（或开球、$\R^n$）；
        \item $O_\alpha\cap O_\beta\ne\mt$ 时 $\psi_\alpha\circ\psi_\beta^{-1}$ 光滑，
    \end{itemize}
    则称 $\{(O,\psi):O\in\topo P\}$ 为光滑图册。若一个光滑图册 $\mathcal A$ 不真含于任何光滑图册，则称 $(M,\topo T,\mathcal A)$ 为一个\textbf{光滑拓扑流形}，简称\textbf{流形}。
\end{definition}
\end{tcolorbox}
或简言之，
\begin{center}
    \tcbox{
    \textbf{流形}是拓扑性质良好的空间，赋有相互光滑变换的一切局部坐标系。
    }
\end{center}

定量计算流形间的映射，必然使用坐标系，即通过 $\psi'\circ f\circ \psi^{-1}$ 来确定 $f$。
\begin{definition}
    置 $C^k$-流形 $M,M'$。$f:M\to M'$ 称为 $C^r$\textbf{-映射}，若 $\forall p\in M$，存在坐标系 $(\psi,O),(\psi',O')$ 使 $p\in O,f(p)\in O'$ 且 $\psi'\circ f\circ \psi^{-1}$ 是 $C^r$ 的（$r\leqslant k$）。
\end{definition}

\begin{remark}
    $r\leqslant k$ 使 $C^r$ 性与坐标选择无关。此后类似概念的良定义性均同理。
\end{remark}

\begin{example}
    坐标变换 $\psi'\circ \psi^{-1}$ 相当于取恒等 $f=\id_M:M\to M$。
\end{example}
\begin{example}
    曲线是流形间 $C^0$-映射 $C:I\to M$ 的特例。其中将 $I\subset\R$ 上坐标系视作恒等映射，从而只需考察\textbf{参数方程} $\psi\circ C$ 的性质。
\end{example}

\begin{definition}
    流形上的曲线为 $C^r$-映射就称 $C^r$\textbf{-曲线}。有限个光滑曲线收尾相连构成\textbf{分段光滑曲线}。
\end{definition}

\begin{definition}
    \textbf{$C^r$-同胚}是流形间双向 $C^r$ 的双射。$C^\infty$-同胚称为\textbf{微分同胚}或\textbf{光滑同胚}。
\end{definition}

\begin{theorem}
    开球和 $\R^n$ 微分同胚
\end{theorem}
\begin{proof}
    不失一般性，$B(0,1),\R^n$ 的同胚式 \eqref{eq:B-Rn} 是双向光滑的。
\end{proof}

微分同胚的流形可建立起等价关系，故我们常考虑的微分同胚实际上是自同胚。
看待微分同胚 $f:M\to M$ 有两种不同观点：将 $f$ 用同一个坐标系表为 $\psi\circ f\circ \psi^{-1}$，此即“点动、坐标不动”的主动观点；定义 $\psi':=\psi\circ f$ 就有坐标变换 $\psi' \circ\psi^{-1}$，此即“点不动、坐标动”的被动观点。

以（光滑拓扑）流形为对象、光滑映射为态射，构成\textbf{流形范畴} $\cate{Man}$，同构态射即微分同胚。


\begin{example}
    流形上的函数是流形间映射 $f:M\to\F$ 的特例。其中将 $\F=\R,\C$ 上坐标系视作恒等映射，从而只需考察 $n$ 元函数 $f\circ\psi^{-1}$ 的性质。
\end{example}

\begin{definition}
    流形上的函数称为\textbf{标量场}(scalar field)。$U\subset M$ 上全体 $C^r$-标量场构成的集合记作 $C^r(U)$。
\end{definition}

我们曾说明 $\F^X$ 是 $\F$-模。原则上还可引入乘法：$\forall p\in X,f,g\in \F^X$，
\eq{(f+g)(p):=f(p)+g(p),\quad  (fg)(p):=f(p)g(p),}
将常函数的赋值 $\alpha\in\F$ 就记为函数名，即 $\alpha(p):=\alpha$，于是数乘即用常函数相乘，同时也就囊括了零元。故 $\F^X$ 同时是 $\F$-模和交换幺环，称\textbf{函数环}。且注意可有 $\alpha(fg)=(\alpha f) g=f(\alpha g)$，从而 $\F^X$ 还是 $\F$-代数。
但 $\F^X$ 不能是域，因为非零元不一定有逆：设 $f\ne 0$ 但存在零点，则不存在 $g$ 使 $f(p)g(p)=1$ 恒成立。

现在考虑 $C^r(U)$。$C^r$-性保持了运算封闭，故 $C^r(U)$ 仍是交换幺 $\F$-代数，可称 $U$ 上的 $C^r$-函数环。为方便，我们常考虑实光滑函数。

\begin{example}
    $C^k$-流形上的第 $i$ \textbf{坐标映射} $x^i:O\to\R$ 属于 $C^k(O)$。
\end{example}

\section{流形的拓扑性质*}

\begin{theorem}
    局部欧 $\im$ 局部路径连通 $\land$ 局部可度量；若再满足 $T_2$，则局部紧。
\end{theorem}
\begin{proof}
    任给 $p\in M$ 及其邻域 $N$，取 $(O,\psi)$ 使 $p\in O$。存在 $\psi(p)$ 为心的欧氏开球 $B\subset \psi[O\cap N^\circ]$，同胚保持路径连通性，则 $\psi^{-1}[B]\subset N$ 是 $p$ 的路径连通开邻域，故 $M$ 局部路径连通。

    坐标域 $O$ 上的度量 $d(p,q):=\|\psi(p)-\psi(q)\|$ 诱导其子拓扑，故 $M$ 局部可度量。

    若 $M$ 为 $T_2$，取上述 $B$ 使闭球 $\bar B\subset\psi[O]$。则 $K:=\psi^{-1}[\bar B]$ 紧，因而在 $M$ 中闭。于是开邻域 $U:=\psi^{-1}[B]$ 在 $M$ 中的闭包满足 $\bar U\subset K$，故紧，即 $M$ 局部紧。
\end{proof}

\begin{theorem}
    路径连通 $\xLeftrightarrow{\text{局部欧}}$ 连通。
\end{theorem}
\begin{proof}
    已知路径连通 $\im$ 连通。反之，局部欧 $\land$ 连通 $\im$ 局部路径连通 $\land$ 连通 $\im$ 路径连通。
\end{proof}

\begin{theorem}
    拓扑流形可度量。
\end{theorem}
\begin{proof}
    定义一：$\text{仿紧} \land T_2 \land \text{局部欧}\im \text{仿紧} \land T_2 \land \text{局部可度量} \xRightarrow{\text{Smirnov}}$ 可度量。

    定义二：$C_2 \land T_2 \land \text{局部欧} \im C_2 \land T_2 \land \text{局部紧} \im C_2 \land T_3 \xRightarrow{\text{Urysohn}}$ 可度量。
\end{proof}
\begin{remark}
    注意可度量 $\im$ 仿紧，所以其实定义二 $\im$ 定义一。
\end{remark}

\begin{definition}
    若开覆盖 $\topo P=\{O_i\}$ 的单位分解 $\{\rho_i\}$ 光滑，则称\textbf{光滑单位分解}。
\end{definition}

\begin{lemma}[光滑 Urysohn 引理]\label{thm:smooth_urysohn_bump}
    设 $M$ 为光滑流形，$K\subset U\subset M$，其中 $K$ 紧、$U$ 开。则存在 $\varphi\in C^\infty(M)$，使 $0\leqslant\varphi\leqslant1$、$\varphi|_K=1$，且 $\supp\varphi$ 为含于 $U$ 的紧集。
\end{lemma}
\begin{proof}
    置 $h(t)=\e^{-1/t}$（$t>0$），$h(t)=0$（$t\leqslant0$）。由前文知 $h$ 光滑，故
    \[
        \chi(t):=\frac{h(1-t)}{h(1-t)+h(t-1/4)}
    \]
    光滑，取值于 $[0,1]$，且在 $t\leqslant1/4$ 时为 $1$，在 $t\geqslant1$ 时为 $0$。

    对每个 $x\in K$，取坐标系 $(O_x,\psi_x)$ 使 $x\in O_x\subset U$、$\psi_x(x)=0$，并取 $r_x>0$ 使 $\bar B(0,r_x)\subset\psi_x[O_x]$。在 $O_x$ 上置
    \[
        b_x(p):=\chi\!\left(\frac{\|\psi_x(p)\|^2}{r_x^2}\right),
    \]
    并在 $M\backslash O_x$ 上置零。其支撑含于紧集 $\psi_x^{-1}[\bar B(0,r_x)]\subset O_x$，故延拓后仍光滑；且 $b_x$ 在开邻域 $V_x:=\psi_x^{-1}[B(0,r_x/2)]$ 上恒为 $1$。

    由紧性，取 $V_{x_1},\ldots,V_{x_N}$ 覆盖 $K$，令
    \[
        \varphi:=1-\prod_{i=1}^N(1-b_{x_i}).
    \]
    则 $\varphi$ 光滑，$0\leqslant\varphi\leqslant1$，在 $K$ 上恒为 $1$，且 $\supp\varphi\subset\bigcup_{i=1}^N\supp b_{x_i}\subset U$，故支撑紧。$K=\mt$ 时取 $\varphi=0$ 即可。
\end{proof}

\begin{theorem}
    光滑流形的任何开覆盖均存在从属于它的局部有限的光滑单位分解。
\end{theorem}
\begin{proof}
    设开覆盖 $\{O_\alpha\}$。由局部紧性，先取相对紧开集组成的加细，再由仿紧性取其局部有限开加细 $\{U_i\}$，其中 $U_i\subset O_{\alpha(i)}$、$\bar U_i$ 紧。沿用连续单位分解证明中的收缩构造，取闭覆盖 $\{F_i\}$ 使 $F_i\subset U_i$，则各 $F_i$ 紧。

    由引理 \ref{thm:smooth_urysohn_bump}，取光滑函数 $b_i$ 使 $0\leqslant b_i\leqslant1$、$b_i|_{F_i}=1$、$\supp b_i\subset U_i$。令
    \[
        b:=\sum_i b_i,\quad
        \rho_\alpha:=\frac{\sum_{\alpha(i)=\alpha}b_i}{b}.
    \]
    局部有限性保证各和光滑，且 $\{F_i\}$ 覆盖 $M$ 保证 $b\geqslant1$。故 $\rho_\alpha$ 光滑，$0\leqslant\rho_\alpha\leqslant1$，$\sum_\alpha\rho_\alpha=1$。又由局部有限性，$\bigcup_{\alpha(i)=\alpha}\supp b_i$ 闭，故
    \[
        \supp\rho_\alpha\subset\bigcup_{\alpha(i)=\alpha}\supp b_i\subset O_\alpha.
    \]
    每点附近只遇到有限个 $U_i$，故也只遇到有限个 $\supp\rho_\alpha$，即得所求。
\end{proof}


\section{切空间}

物理学的矢量总是靠作差或求导来产生，即切矢。曾指出给定某几何形体的参数方程 $x^\mu(t,\cdots)$ 后，可得到切矢场 $T^\mu = \pdv*{x^\mu}{t}$，它具有变换 $T'^\nu = \pdv{x'^\nu}{x^\mu}T^\mu$。而其中 $\pdv*{x'^\nu}{x^\mu}$ 就是 $\{x\}$ 的坐标基，故这相当于将切矢展开到 $\{x\}$ 坐标基 $T'^\nu = T^\mu \pdv{x'^\nu}{x^\mu}$，但最终用 $\{x'\}$ 表达。
可见，坐标分量 $T'^\nu$ 是 $\pdv{t}=\pdv{x^\mu}{t}\pdv{x^\mu}=T^\mu \p_\mu$ 对 $x'^\nu:O\to\R$ 作用的结果，故仅需偏导数本身配合所作用的 $x'^\nu(x)$ 即可给出我们对切矢要求的所有性质。
立即发现，这种描述已独立于坐标语言。数学又常常使用映射来构造概念，故何不直接定义 $T=\pdv*{t}$？
Cartan 最先提出这一思想：用一个具有求导性质的泛函来代指切矢。

将切矢定义为一个求导映射，与最传统的曲线相切有何直观联系？
给定一条穿过 $p\in M$ 的曲线 $\gamma$ 且在 $p=\gamma(0)$ 处可导，规定给出相同切矢的曲线等价，相同性通过坐标分量 $\dv{t}\big|_{t=0}(\psi\circ\gamma)$ 来判断。
置任意可导标量场 $f$，若某点切矢为 $A^\mu$，则方向导数为 $\dv*{f}{t}=A^\mu\p_\mu f$。故等价类 $[\gamma]$ 中的曲线给出了相同的方向导数 $\dv*{t}$，即曲线等价类与方向导数一一对应。
标量场在此处的作用是引出切矢。犹如在电场中放入测试电荷，用静电力研究电场一样，标量场受导数作用出实数，以研究导数自身，而标量场不是关键。

剩下的问题在于需要哪些求导性质。若直接使用偏导符号 $A=A^\mu\pdv{x^\mu}$，实际上仍然使用了坐标语言，因此要将目光转向导数的代数性质。
若 $f$ 是解析的，展开为
\[
    f=f(p)+\p_i f(p)(x^i-x^i(p))+ \frac 12\p_i\p_j f(p)(x^i-x^i(p))(x^j-x^j(p))+\cdots.
\]
$A=A^\mu\p_\mu$ 是求导操作，至少要考虑到线性性、Leibniz 乘法法则 $A(fg)=g(p)A(f)+f(p)A(g)$。
乘法法则给出 $A(1)=A(1\cdot 1)=2A(1)\Rightarrow A(1)=0$，再由线性性知$\forall c\in\R$，$A(c)=A(c\cdot 1)=c A(1)=0$。
从而 $A$ 作用于 $f$ 的确只剩下 $A(f)=A(x^\mu)\p_\mu f(p)$，高阶项为零是因为必然存在 $x^j(p)-x^j(p)=0$。可见线性性、乘法法则已能反求得 $A=A^\mu\p_\mu$，故此即所需的全部求导代数性质了。
当然，对无穷级数滥用代数运算是危险的，其次也不需要 $f$ 解析，我们更倾向于考虑有限项相加。

需要注意，坐标函数 $x^\mu$ 通常只在坐标域上定义，好在求导只关心 $p$ 附近的函数值：只要 $M$ 上函数 $f_1,f_2$ 在 $U\in\topo N(p)$ 上相等，置 $f:=f_1-f_2$，令 $h$ 在 $M\backslash U$ 上为零但 $h(p)\ne 0$，就可构造零元 $fh=0$，则 $0=A(fh)=h(p)A(f)$ 从而 $A(f)=0$，即 $p$ 点任意切矢 $A$ 满足 $A(f_1)=A(f_2)$。便能合法谈及坐标分量 $A(x^\mu)$。

\begin{definition}
    置光滑流形 $M$ 及 $p\in M$。设 $U,V$ 为 $p$ 的开邻域，$f\in C^\infty(U),g\in C^\infty(V)$ 为实函数。若存在开邻域 $W\ni p$，使 $W\subset U\cap V$ 且 $f|_W=g|_W$，则称 $f,g$ 在 $p$ 处等价。等价类 $[f]_p$ 称为 $p$ 处的\textbf{函数芽}，全体记作 $C^\infty_p(M)$。在共同定义域上作运算，规定
    \eq{
        [f]_p+[g]_p:=[f+g]_p,\quad c[f]_p:=[cf]_p,\quad [f]_p[g]_p:=[fg]_p,
    }
    则 $C^\infty_p(M)$ 构成实代数，且 $[f]_p(p):=f(p)$ 与代表元无关。

    $p$ 处的\textbf{切矢}是满足 Leibniz 乘法法则的实线性映射 $v:C^\infty_p(M)\to\R$，即
    \eq{
        v([f]_p[g]_p)=g(p)v([f]_p)+f(p)v([g]_p).
    }
    对任何在 $p$ 附近定义的光滑函数，总可简记 $v([f]_p)=v(f)$。
    $p$ 处全体切矢的集合 $T_p M$ 称为 $M$ 在 $p$ 处的\textbf{切空间}。全体切空间的无交并
    \eq{
        TM:=\bigsqcup_{p\in M} T_p M=\bigcup_{p\in M}\{p\}\times T_p M=\{(p,v):p\in M,v\in T_p M\}
    }
    称为\textbf{切丛}(tangent bundle)。
\end{definition}

\begin{remark}
    直观图像是把切空间理解为背景在该点的切平面，切丛就是全体切平面。

    此处直接考虑光滑性 $r=\infty$，是为了使坐标分解
    \[
        f-f(p)=\sum_i(x^i-x^i(p))H_i
    \]
    中的 $H_i$ 仍属于同一函数代数，从而可对各项使用 Leibniz 法则。对 $f\in C^r$（$1\leqslant r<\infty$），这种分解的系数一般仅为 $C^{r-1}$，$v(H_i)$ 未必有定义；光滑时则无此问题，并可由此证明切矢恰由 $\dim M$ 个坐标分量决定。因此不能直接将上述定义中的 $\infty$ 换为有限 $r$；有限阶可微流形可用曲线等价类定义切矢。
\end{remark}

\begin{theorem}
    前文已暗示更简单的等价定义：$p$ 处的切矢是满足
    \[
        D(fg)=g(p)D(f)+f(p)D(g)
    \]
    的实线性泛函 $D:C^\infty(M)\to\R$。两种定义的对应关系为 $D(f)=v([f]_p)$。
\end{theorem}
\begin{proof}
    由函数芽上的切矢 $v$ 定义 $D(f):=v([f]_p)$，显然满足上述条件。

    反之，先证 $D$ 只依赖函数芽。若 $h\in C^\infty(M)$ 在 $p$ 的开邻域 $W$ 上恒为零，取光滑函数 $\chi$ 使 $\chi(p)=1$、$\supp\chi\subset W$，则 $h\chi=0$，故
    \[
        0=D(h\chi)=\chi(p)D(h)+h(p)D(\chi)=D(h).
    \]
    因而在 $p$ 附近相等的全局函数具有相同的 $D$ 值。

    任给函数芽 $[f]_p$，设 $f\in C^\infty(U)$。取光滑函数 $\rho$，使其在 $p$ 的某邻域上恒为 $1$，且 $\supp\rho\subset U$。将 $\rho f$ 在 $M\backslash U$ 上置零，得到 $F\in C^\infty(M)$，且 $[F]_p=[f]_p$。于是 $v([f]_p):=D(F)$ 与代表元及延拓的选择无关，并保持线性性与 Leibniz 法则。每个芽都有这样的全局代表，故此对应唯一且互逆。
\end{proof}

\begin{theorem}
    $T_p M$ 是线性空间，且 $\dim T_p M=\dim M =: n$。
\end{theorem}
\begin{proof}
    切矢对通常的泛函加法及数乘封闭，故 $T_p M$ 为实线性空间。取坐标系 $(O,\psi)$，置 $a:=\psi(p)$。由求导法则，
    \[
        \p_i|_p(f):=\pdv{(f\circ\psi^{-1})}{x^i}\bigg|_a
    \]
    是切矢，且 $\p_i|_p(x^j)=\delta_i^j$，故这 $n$ 个切矢线性无关。此即 $\dim T_p M\geqslant n$。

    任给 $f$ 在 $p$ 处的光滑函数芽，缩小 $O$，使 $f$ 在 $O$ 上光滑且 $\psi[O]$ 为以 $a$ 为心的开球。记 $\hat f:=f\circ\psi^{-1}$，由微积分基本定理（对 $i$ 用求和约定），
    \[
        f(q)-f(p)= (x^i(q)-x^i(p))H_i(q),\quad
        H_i(q):=\int_0^1\d{t}\p_i\hat f\bigl(a+t(\psi(q)-a)\bigr).
    \]
    各 $H_i$ 光滑，且 $H_i(p)=\p_i|_p(f)$。任给 $v\in T_p M$，由 $v(c)=0$ 及 Leibniz 法则，得
    \[
        v(f)= v(x^i)H_i(p)= v(x^i)\p_i|_p(f).
    \]
    因 $f$ 任意，故
    \eq{\label{chdm:eqn_v-expand-on-natural-bases}
        v= v^i\p_i|_p,\quad v^i:=v(x^i).
    }
    于是 $\{\p_i|_p\}_{i=1}^n$ 张成 $T_p M$，从而是一组基，故 $\dim T_p M=n$。
\end{proof}

\begin{definition}
    $X_i = \p_i|_p$ 称为 $p$ 处第 $i$ 个\textbf{坐标基矢}。$\{X_i\}$ 称为 $p$ 处的\textbf{坐标基}。
\end{definition}

\begin{theorem}[坐标分量变换律]
    置坐标系 $(O,\psi),(O',\psi')$ 及 $p\in O\cap O'$。两坐标基及同一切矢 $v\in T_pM$ 的分量满足
    \[
        X_i= {\pdv{x'^j}{x^i}}(p)X'_j,\quad
        v'^j= {\pdv{x'^j}{x^i}}(p)v^i, \quad {\pdv{x'^j}{x^i}}(p):=\pdv{(x'^j\circ\psi^{-1})}{x^i}\bigg|_{\psi(p)}.
    \]
    其中 Jacobi 系数 $\pdv*{x'^j}{x^i}$ 是 $O\cap O'$ 上的光滑函数。
\end{theorem}
\begin{proof}
    任给 $p$ 处的光滑函数芽 $f$，由 $f\circ\psi^{-1}=(f\circ\psi'^{-1})\circ(\psi'\circ\psi^{-1})$ 及链式法则，得
    \[
        X_i(f)= {\pdv{x'^j}{x^i}}(p)X'_j(f).
    \]
    因 $f$ 任意，即得坐标基变换。将其代入 $v= v^iX_i= v'^jX'_j$，比较系数即得分量变换。
\end{proof}

切矢场就是给每点指定该点的一个切矢：
\begin{definition}
    置开集 $U\subset M$。若映射 $v:U\to TM$ 满足 $v(p)\in T_pM$（$\forall p\in U$），则称其为 $U$ 上的\textbf{切矢场}，亦称 $TM$ 在 $U$ 上的\textbf{截面}(section)。

    若 $v=v^i\p_i$ 在各坐标域中的分量 $v^i$ 均为 $C^r$ 函数，则称 $v$ 为 $C^r$ 切矢场。全体光滑切矢场记作 $\Gamma(TM|_U)$，亦记 $\mathfrak X(U)$，其中 $TM|_U:=\bigsqcup_{p\in U}T_pM$。$U=M$ 时简记 $\Gamma(TM)$。本书不标可微阶数的 $\Gamma$ 均指光滑截面的集合。
\end{definition}

\begin{definition}
    置 $C^1$-曲线 $\gamma:I\to M$。其在 $p=\gamma(t)$ 处的\textbf{切矢} $\dot\gamma(t)\in T_pM$ 定义为
    \eq{\label{chdm:eqn_DiffCurve}
        \dot\gamma(t)(f):=\dv{(f\circ\gamma)}{t},\quad f\in C^\infty_p(M).
    }
    由链式法则，其坐标展开为 $\dot\gamma(t)=\dv{(x^i\circ\gamma)}{t}\p_i|_{\gamma(t)}$，时常简记 
    \eq{\label{chdm:eqn_DiffCurve-2}
        \dv{t} = \dv{x^i}{t}\p_i.
    }
    映射 $t\mapsto\dot\gamma(t)$ 是沿 $\gamma$ 的\textbf{曲线切矢场}。曲线在 $p$ 处的切矢属于 $T_pM$，而 $T_pM$ 中的切矢未必都切于该曲线。
\end{definition}

\begin{example}
    第 $x^i$ 坐标线的切矢正是坐标基 $\delta^j_i\p_j=\p_i$。
\end{example}


\begin{definition}
    置映射 $f,g:A\to A$。若 $f\circ g=g\circ f$，则称 $f,g$ \textbf{对易}(commute)。若 $A$ 为加法交换群，则可逐点定义其\textbf{对易子}(commutator)
    \[
        [f,g](x):=f(g(x))-g(f(x)).
    \]
    此时 $f,g$ 对易 $\iff[f,g]=0$，其中 $0$ 表示零映射。
\end{definition}

\begin{remark}
    此定义不要求 $f,g$ 线性，但一般映射的对易子未必具有双线性或满足 Jacobi 恒等式。在实线性空间的线性自映射上，对易子具有上述性质，构成通常的 Lie 括号。
\end{remark}

光滑切矢场 $A$ 可视为算子 $C^\infty(U)\to C^\infty(U)$，规定 $A(f)(p):=A_p(f)$。以下 $AB$ 指这种算子的复合，即 $AB(f)=A(B(f))$。

\begin{example}
    微积分 Schwarz 定理指出，光滑函数的混合偏导可交换：$ [\p_\mu,\p_\nu ]=0$，称为\textbf{偏导对易性}。
\end{example}

\begin{theorem}
    光滑矢量场对易子 $[A,B]$ 仍是光滑矢量场。
\end{theorem}
\begin{proof}
    $[A,B]$ 显然线性。展开 $A(B(fg))-B(A(fg))$，交叉项抵消，得
    \[
        [A,B](fg)=g[A,B](f)+f[A,B](g).
    \]
    故在每点取值均为切矢。又在坐标域中，二阶偏导项相消，得
    \eq{
        [A,B]^\lambda=A^\mu\partial_\mu B^\lambda-B^\mu\partial_\mu A^\lambda.
    }
    各分量均光滑，故 $[A,B]$ 为光滑切矢场。如果只使用分量语言，还需额外再证它满足张量变换式：
    \[
        \pdv{x^\sigma}(\pdv{x'^\lambda}{x^\kappa}B^\kappa)=  B^\kappa\pdv{x'^\lambda}{x^\sigma}{x^\kappa}+\pdv{x'^\lambda}{x^\kappa}\partial_\sigma B^\kappa,
    \]
    因为 $\partial_\sigma\partial_\kappa=\partial_\kappa\partial_\sigma$ 始终使 $\pdv{x'^\lambda}{x^\sigma}{x^\kappa}$ 项抵消。
\end{proof}

\begin{definition}
    $T^*_p M:=(T_p M)^*$ 称为 $M$ 在 $p$ 处的\textbf{余切空间}，其元素称为 $M$ 在 $p$ 处的\textbf{余切矢}。
    全体余切空间的无交并 $T^*M:=\bigsqcup_{p\in M} T^*_p M$ 称为\textbf{余切丛}。
\end{definition}

Cartan 不仅提出了切矢的泛函导子定义，他还注意到
%% ...
这相当于给函数微分一个更严格的理解。标准分析学中常将 $\d f$ 定义为 $\Delta f$ 的线性主部，即 $\d f(x,\Delta x):=f'(x)\Delta x$。这里 $\Delta x$ 是从某点 $x$ 出发的位移，实际上任意大小，但物理上的理解是在 $\Delta x$ 很小时给出的微小增量。
%% ...
可见微分和余切的关系。

\begin{theorem}
    $\dim T_p^* M = \dim T_p M = n$。
\end{theorem}
\begin{proof}
    $\{\p_i|_{p}\}$ 的对偶基记为 $\{\left.\d x^j\right|_{p}\}$，由
\eq{\label{chdm:eqn_dual-natural-bases}
    \left.\d x^j\right|_{p}\!\left(\p_i|_{p}\right):=\delta_i^j
}
确定。下文略去取值点，则 $\d x^j(v)=v^i\delta_i^j=v^j$，即取出 $v$ 的第 $j$ 个分量。任给 $\omega\in T^*_{p}M$，有
\eq{
    \omega=\omega_j\d x^j,\quad \omega(v)=\omega_jv^j.
}
\end{proof}

配对 $\langle\omega,v\rangle:=\omega(v)$ 是实数，而 $v(f)$ 也是实数，因此可用函数 $f$ 定义余切矢量。

\begin{definition}\label{chdm:def_pushforward-vector}
    映射 $F:M\to N$ 的在 $p$ 处的切映射
    定义\textbf{推前映射}，或称为\textbf{切映射}：
    \eq{
        (\phi_{*}v)(g) := v(\phi^{*}(g)) = v (g \circ \phi),
        \quad \forall g\in C^\infty_{\phi(p)}(N),\quad \forall v \in T_{p}M .
    }
\end{definition}

\begin{definition}函数的\textbf{微分}定义为
    \eq{
    \d f(p)(A):=A(f).
}
\end{definition}

由定义及坐标展开，
\eq{\label{chdm:eqn_dfv}
    (\d f)_{p}(v)=v(f)=v^i\left.\pdv{f}{x^i}\right|_{p}.
}
特别地，$\d x^i(\p_j)=\p_j(x^i)=\delta^i_j$，故坐标函数的微分正是上述对偶基，称为\textbf{对偶坐标基}。于是
\eq{\label{chdm:eqn_df}
    (\d f)_{p}=\left.\pdv{f}{x^j}\right|_{p}\left.\d x^j\right|_{p}.
}
略去取值点，即 $\d f=\p_jf\,\d x^j$，与通常的全微分一致。

在标准分析中，全微分理解为各线性映射间的关系，这里将线性映射叙述为余切矢量。由于分量是 $\p_i f$，因此有时建议称标量场的\textbf{梯度}\footnote{为区分于标准分析学中的概念，部分书籍会用粗体表示这些泛函，如 $\bm A,\bm\omega,\bm{\partial}/{\bm{\partial}x^i},\mathbf{d}f$。但其实没有必要，因为两种概念其实是一致的。看似不同是因为本科教学为了通俗性，默认以坐标空间为基础，而没有额外提出流形。}。

函数的微分，对于标量函数有这些运算：加法 $f+g$、乘法 $fg$ 和微分 $f\mapsto\d f$。乘法对于加法是可分配的。微分对乘法也是一个分配，即 Leibniz 法则。

对于余切矢量认知，也可以先把式\eqref{chdm:eqn_dfv}当成定义，然后由此
立刻导出坐标基矢间的内积公式\eqref{chdm:eqn_dual-natural-bases}；再导出其它公式。
很明显定义\eqref{chdm:eqn_dual-natural-bases}和式\eqref{chdm:eqn_dfv}是相容的。


微分流形定义中要求坐标卡可以取成各种相容形式，如果$p\in M$有两个相互容许的
坐标卡$(U,\varphi;x)$和$(V,\psi;y)$，且$U\cap V \neq \varnothing$；
对应的自然坐标基底分别是$(\frac{\partial }{\partial x^i})_{p} $
和$(\frac{\partial }{\partial y^j})_{p} $。
我们把$(\frac{\partial }{\partial y^j})_{p} $看成
式\eqref{chdm:eqn_v-expand-on-natural-bases}中的$v$，那么有
\eq{\label{chdm:eqn_xy-transform}
    \left(\frac{\partial }{\partial y^j}\right)_{p} = \left. \frac{\partial x^i}{\partial y^j}\right|_{p} \cdot
    \left(\frac{\partial }{\partial x^i}\right)_{p}     {\ \Leftrightarrow \ }
    \left(\frac{\partial }{\partial x^i}\right)_{p} = \left. \frac{\partial y^j}{\partial x^i}\right|_{p} \cdot
    \left(\frac{\partial }{\partial y^j}\right)_{p} .
}
上式同时写出了逆变换；可见过渡矩阵恰好是$p$点的Jacobi矩阵。


在$U\cap V$中可以把$y^j$看成自变量是$x^i$的标量函数（也可反过来），
将式\eqref{chdm:eqn_df}中的$f$用$y^j$替换后，有
\eq{\label{chdm:eqn_xy-cov-transform}
    (\d y^j) = (\d x^i) \left.\frac{\partial y^j}{\partial x^i}\right|_{p}
    {\quad \Leftrightarrow \quad }
    (\d x^i) = (\d y^j) \left.\frac{\partial x^i}{\partial y^j}\right|_{p}
}
过渡矩阵恰好是$p$点的Jacobi矩阵之逆。

其实，在同一流形$M$中，点$p\in M$有两个相互容许的坐标卡$(U,\varphi;x)$和$(V,\psi;y)$，
且$U\cap V \neq \varnothing$；此时的局部坐标变换与局部微分同胚是等价的，
一般分别称它们为被动模式和主动模式。请读者仔细理解这两种操作的等价性。

有了切空间$T_{p}M$及其对偶空间$T_{p}^{*}M$（两个都是线性空间），
仿照\S\ref{sec:tensor}，可以定义微分流形$M$在点$p$处的$({p},{q})$型张量，
可以把它看成$m^{p+q}$维线性空间
\eq{\label{chdm:eqn_TM-TsM}
    \mathcal{T}^p_q(p) := \underbrace{T_{p}M\otimes \cdots \otimes T_{p}M}_{p\text{个}}
    \otimes \underbrace{T_{p}^{*}M\otimes \cdots \otimes T_{p}^{*}M}_{q\text{个}}
}
中的元素。也可以把它看成笛卡尔积空间
\eq{\label{chdm:eqn_TM-X-TsM}
    \underbrace{T_{p}^{*}M \times \cdots \times T_{p}^{*}M}_{p\text{个}}
    \times \underbrace{T_{p}M\times \cdots \times T_{p}M}_{q\text{个}}
}
的$p+q$重线性函数。


从上一小节可知线性空间$\mathcal{T}^p_q(p)$的自然基矢量是（略去了下标$p$）
\eq{\label{chdm:eqn_tensortrans}
    \p_{i_1}\otimes\cdots\otimes\p_{i_p}\otimes
    \d x^{j_1}\otimes\cdots\otimes\d x^{j_q}.
}
上式中$1\leqslant \{i,j\} \leqslant m $。
由式\eqref{chdm:eqn_xy-transform}和\eqref{chdm:eqn_xy-cov-transform}不难得到
上式在坐标变化下的变换关系。

任一张量 $\alpha\in\mathcal T^p_q(p)$ 均可展开为
\eq{\label{chdm:eqn_tensor-component}
    \alpha=\alpha^{i_1\cdots i_p}{}_{j_1\cdots j_q}
    \p_{i_1}\otimes\cdots\otimes\p_{i_p}\otimes
    \d x^{j_1}\otimes\cdots\otimes\d x^{j_q}.
}

由式\eqref{chdm:eqn_xy-transform}及\eqref{chdm:eqn_xy-cov-transform}，坐标变换时的分量满足
\eq{\label{chdm:eqn_tensor-component-trans}
    \tilde{\alpha}^{k_1 \cdots k_p}_{\phantom{k_1 \cdots k_p}l_1\cdots l_q} (y)
    ={\alpha}^{i_1 \cdots i_p}_{\phantom{i_1 \cdots i_p}j_1\cdots j_q} (x) \times
    \frac{\partial y^{k_1}}{\partial x^{i_1}} \cdots
    \frac{\partial y^{k_p}}{\partial x^{i_p}} \times
    \frac{\partial x^{j_1}}{\partial y^{l_1}} \cdots
    \frac{\partial x^{j_q}}{\partial y^{l_q}} .
}
需要强调的是：上式所有量都需在点$p\in M$取值，
$\{x\}$、$\{y\}$是$p$点某个小邻域的两套坐标系。
由此式直接可得一阶切矢量或余切矢量的分量变换关系。

\section{诱导映射}\label{chdm:sec_induced-map}

考虑对偶函子 $D$。

设有两个微分流形$M$和$N$，维数分别为$m$和$n$（两者未必相等），存在光滑映射$\phi:M\to N$；
$M$和$N$的局部坐标分别为$(U;x^i)$和$(V;y^\alpha)$，则
\eq{
    y^\alpha = \phi(x^i)=y^\alpha (x^1,\cdots, x^m), \quad 1\leqslant \alpha \leqslant n .
}
取微分流形$M$中任一点$p$及微分流形$N$中点$\phi(p)$，
则有$y^\alpha_0=\phi(x^i_0)$。
那么它自然诱导出一系列映射。


给定标量场 $f:N\to\R$，$p\in M$ 变到 $\phi(p)\in N$ 后标量场的值为 $f(\phi(p))$，而我们可将 $\phi(p)$ 处的 $f$ 值直接搬运回 $p$ 处，即定义新标量场 $\phi^* f:M\to\R$ 满足 $\phi^* f(p):=f(\phi(p))$。
考虑拉回 $\phi^*:C^\infty(N)\to C^\infty(M)$，即 $\phi^* f:=f\circ\phi$。
易知 $\phi^*$ 是线性映射，且 $\phi^*(fg)=(\phi^* f) (\phi^* g)$。
类似地，记 $\phi(p)=q$，若 $\phi$ 是同胚，则也可将 $p=\phi^{-1}(q)$ 处的 $f$ 值往前搬至 $q$ 处，即定义 $\phi_* f$ 使 $\phi_* f(q):=f(\phi^{-1}(q))$ 或者 $\phi_* f:=f\circ\phi^{-1}$

\begin{definition}\label{chdm:def_pullback-scalar}
    从$C^\infty_{\phi(p)}(N)$到$C^\infty_{p}(M)$
    的拉回 $\phi^{*}$ 为
    \eq{
        \phi^{*}(g) := g \circ \phi,
        \quad \forall g\in C^\infty_{\phi(p)}(N) .
    }
\end{definition}
一般来说，我们不能把$M$中的标量函数推前到$N$上；比如映射$\phi$不是单射，
假设有$\phi(x_1)=\phi(x_2)=y\in N$，那么$y$处的被推前函数值就不知道该
选为$g(x_1)$还是$g(x_2)$的值了，出现多值选择困难。


同样注意，对于一般的光滑映射，我们不能定义切矢量的拉回；请读者
仿照上面讨论，用非单一性给出（类似）解释。

\begin{definition}\label{chdm:def_pullback-1form}
    \textbf{拉回映射}可拓展到余切矢量，也称为\textbf{余切映射}：
    \eq{
        (\phi^*\omega)(v):=\omega(\phi_*v) ,
        \quad \forall \omega \in T^{*}_{\phi(p)}N,\quad \forall v \in T_{p}M .
    }
\end{definition}
需要强调的是：上面几个定义都是点对点的。


\begin{proposition}
    切映射和余切映射都是线性的。
\end{proposition}
\begin{proof}
    先证切映射是线性的。设$\forall u,v \in T_{p}M,\  \forall g\in C^\infty_{\phi(p)}(N),\
    \forall \lambda \in \R$，

    \begin{align*}
        \bigl(\phi_{*}(u+ \lambda v) \bigr)(g)= (u+ \lambda v)(g\circ \phi)
        =u(g\circ \phi) +\lambda \cdot v (g\circ \phi)
        =(\phi_{*} u) (g) + \lambda \cdot (\phi_{*}v)(g) .
    \end{align*}

    再证余切映射的线性性。任给 $\omega,\mu\in T^*_{\phi(p)}N$、$v\in T_{p}M$ 及 $\lambda\in\R$，有
    \[
        \bigl(\phi^*(\omega+\lambda\mu)\bigr)(v)
        =(\omega+\lambda\mu)(\phi_*v)
        =(\phi^*\omega)(v)+\lambda(\phi^*\mu)(v).
    \]

    标量函数的拉回映射\ref{chdm:def_pullback-scalar}也是线性的，
    证明过程完全相同，略。
\end{proof}

\begin{proposition}
    证明切映射确实是“切矢量”，即验证$\phi_{*}v$满足定义\ref{chdm:def_tang-vec-inManifold}。
\end{proposition}
\begin{proof}
    留给读者当作习题。
\end{proof}

\begin{proposition}\label{chdm:thm_coortrans}
    导出自然基矢在切映射和余切映射下的变换关系。
\end{proposition}
\begin{proof}
符号同本小节开头叙述。命题要求解的是切映射$\phi_{*}(\frac{\partial}{\partial x^i})$在局部坐标$(V;y)$的
自然基矢$(\frac{\partial}{\partial y^\alpha})$下的展开系数，
我们只需将它作用在局部坐标$y^\alpha(x^i)$上即可。

\eq{\label{chdm:eqn_push-bases}
    \left[\phi_{*}\left(\frac{\partial}{\partial x^i}\right)\right] (y^\alpha)
    =\frac{\partial}{\partial x^i} \left(y^\alpha \circ \phi \right)
    =\frac{\partial y^\alpha}{\partial x^i}
    \ \Leftrightarrow\ \phi_{*}
    \left(\frac{\partial}{\partial x^i}\right) = \frac{\partial y^\alpha}{\partial x^i}
    \left(\frac{\partial}{\partial y^\alpha}\right)
}
切映射在局部坐标基矢间的关系是Jacobi矩阵，
注意这个矩阵未必是方阵。

余切映射作用于对偶基时，由定义得
\[
    (\phi^*\d y^\alpha)(\p_i)
    =\d y^\alpha(\phi_*\p_i)
    =\pdv{(y^\alpha\circ\phi)}{x^i},
\]
故
\eq{\label{chdm:eqn_pull-bases}
    \phi^*(\d y^\alpha)=\pdv{(y^\alpha\circ\phi)}{x^i}\d x^i.
}

余切映射局部坐标基矢间的关系也是Jacobi矩阵。

最后强调一点：所有公式只在$p$点取值，为使下标简洁上面表达式中略去了此点。
\end{proof}

命题\ref{chdm:thm_coortrans}中映射是$\phi:M\to N$，如果我们令$N=M$，并且
令$\phi$是恒等映射$\id$；但$M$和$N$的局部坐标未必相同。
这时上面命题中的变换可以理解成同一局部邻域内的{坐标变换}；
公式与\S\ref{chdm:sec_coortrans}中的相同。

利用式\eqref{chdm:eqn_pull-bases}可证明如下一个重要公式，
即普通微分和拉回映射可对易。
\begin{align}
    {\phi ^*} (\d f) =& {\phi ^*} \left[ \frac{ \partial f  }
      {\partial y^\alpha} (\d y^\alpha)  \right]
    = \left. \frac{ \partial f }{\partial y^\alpha} \right|_{\phi(p)}
      \left. \frac{{\partial {y^\alpha }}}{{\partial {x^i}}} \right|_{p}
      \left( {\d {x^i }} \right)
    = \left. \frac{ \partial (f \circ \phi) }{\partial x^i} \right|_{p}
      \left( {\d {x^i }} \right)  \notag\\
    =& \d (f \circ \phi) =\d ({\phi ^*}f) . \label{chdm:eqn_dfphi}
\end{align}
仿照上式，其实式\eqref{chdm:eqn_pull-bases}也
可以写为${\phi ^*} (\d y^\alpha) = \d (y^\alpha \circ \phi)$。

给出复合映射的关系式。
\begin{theorem}\label{chdm:thm_map-chain-rule}
    设有光滑流形$M,N$和$Q$，存在流形间的光滑映射$\phi:M\to N$和$\psi:N\to Q$，则
    复合映射$ \psi\circ\phi  : M\to Q$的三种诱导映射服从如下链式法则（其中$\forall p\in M$），

    (1) $(\psi\circ \phi)^{*}= \phi^{*}\circ \psi^{*}:  C^{\infty}_{\psi\circ\phi(p)}(Q)\to C^{\infty}_p (M)$；
    对应定义\ref{chdm:def_pullback-scalar}；

    (2) $(\psi\circ \phi)_{*}= \psi_{*}\circ \phi_{*}: T_p M\to T_{\psi\circ\phi(p)}Q$；
    对应定义\ref{chdm:def_pushforward-vector}；

    (3) $(\psi\circ \phi)^{*}= \phi^{*}\circ \psi^{*}:  T^{*}_{\psi\circ\phi(p)}Q\to T^{*}_p M$；
    对应定义\ref{chdm:def_pullback-1form}。
\end{theorem}
\begin{proof}
    设有任意$v\in T_pM$和任意$f\in C^\infty_{\psi\circ\phi(p)}(Q)$，
    以及任意的$\omega\in T^{*}_{\psi\circ\phi(p)}(Q)$。
    先证第(1)式，由定义直接计算，得
    \begin{align*}
        (\psi\circ \phi)^{*} (f) &= f\circ \psi\circ \phi (p) = (f\circ \psi )\circ \phi (p)
        =\bigl(\psi^*(f) \bigr)\circ \phi (p) = \left. \phi^*\bigl(\psi^*(f) \bigr) \right|_p;
        \ \text{或者} \\
        (\psi\circ \phi)^{*} (f) &= f\circ \psi\circ \phi (p) = (f\circ \psi )\circ (\phi)|_p
        =\phi^* (f\circ \psi )|_p = \left. \phi^*\bigl(\psi^*(f) \bigr) \right|_p .
    \end{align*}
    上式最后一式便是第(1)条中的结论。再证第(2)式：
    \begin{equation*}
        (\psi\circ \phi)_{*}v (f) = v \bigl((\psi\circ \phi)^{*} (f)\bigr)
          =v \Bigl(\phi^{*}\circ \bigl( \psi^{*} (f)\bigr) \Bigr)
          =\phi_{*} v \bigl( \psi^{*} (f)\bigr) =\psi_{*} \circ \phi_{*} v (f) .
    \end{equation*}
    证明过程中用到了第(1)式。最后证明第(3)式：
    \begin{equation*}
        \bigl((\psi\circ \phi)^{*} \omega\bigr) (v) = \omega \bigl((\psi\circ \phi)_{*} v\bigr)
        = \omega (\psi_{*} \circ \phi_{*}  v )
        = \psi^{*} \omega ( \phi_{*}  v)
        = (\phi^{*}\circ \psi^{*} \omega) (v) .
    \end{equation*}
    需要注意，第(1)、(3)式的链式法则次序与原复合映射次序相反。
\end{proof}

\begin{example}\label{chdm:exm_phiD=Dphi}
    曲线$\gamma(t)$切矢的推前等于“曲线像$\phi\circ\gamma(t)$”的切矢。
\end{example}
设有光滑流形$M$和$N$，存在流形间的光滑映射$\phi:M\to N$。
$M$中有光滑参数曲线$\gamma(t)$，此线切矢量如式\eqref{chdm:eqn_DiffCurve-2}所述。
很明显$\phi\circ\gamma:(-\delta,\delta)\to N$是流形$N$中一条光滑曲线，下面来给出
此条曲线的切矢量，$\forall f\in C^\infty_{\phi(0)}(N)$，
\begin{align*}
    \phi_{*}\left(\left.\frac{\d }{\d t}\right|_{\gamma(0)}\right) (f)
    = \left.\frac{\d }{\d t}\right|_{\gamma(0)} (f\circ \phi)
    = \left.\frac{\d }{\d t}\right|_{t=0} \bigl(f\circ \phi\circ\gamma(t)\bigr)
    = \left.\frac{\d }{\d t}\right|_{\phi\circ\gamma(0)} (f) .
\end{align*}
从上式可以得到
\eq{\label{chdm:eqn_phiD=Dphi}
    \phi_{*}\left(\left.\frac{\d }{\d t}\right|_{\gamma(0)}\right)
    = \left(\left.\frac{\d }{\d t}\right|_{\phi\circ\gamma(0)}\right).
}
在涉及曲线切矢量的计算时，此式非常有用。


\begin{exercise}\label{chdm:exer_fg}
	有映射$f(t)=(t,t^2)$，$g(x,y)=x \cos y$。
	求切映射：$f_*$、$g_*$、$(f\circ g)_*$。
\end{exercise}

\begin{exercise}
	算出例题\ref{chdm:exer_fg}中曲线$f$的切矢量，然后推前到像集中；
	先将曲线$f$推前到像集中，再求它的切矢量；进而验证式\eqref{chdm:eqn_phiD=Dphi}是否正确。
\end{exercise}

\section{嵌入映射}\label{chdm:sec_sub-manifold}


\begin{theorem}\label{chdm:thm_inv-in-M}
    设光滑流形$M$和$N$的维数都是$m$，$f:M\to N$是光滑映射。
    如果在一点$p\in M$，切映射$f_{*}:T_p M\to T_{f(p)}N$是同构的，
    则存在点$p$在$M$中的邻域$U$使得$V=f(U)$是点$f(p)$在$N$中的一个邻域，
    并且$f|_U :U\to V$是微分同胚映射。
\end{theorem}
利用上面定理可证明如下定理。
\begin{theorem}\label{chdm:thm_immerse-1}
    设光滑流形$M$和$N$的维数分别是$m$和$n$，且$m<n$。
    $f:M\to N$是光滑映射，并且切映射$f_{*}:T_p M\to T_{f(p)}N$是非退化的。
    则存在点$p$的局部坐标系$(U;x^i)$以及点$q=f(p)$的局部坐标系$(V;y^\alpha)$，
    使得$f(U)\subset V$，并且映射$f|_U$可用局部坐标表示为
    \eq{
        \begin{cases}
            y^i \circ f = x^i, & 1 \leqslant i \leqslant m ; \\
            y^\nu \circ f = 0, & 1+m \leqslant \nu \leqslant n .
        \end{cases}
    }
\end{theorem}
\begin{proof}
    设光滑映射$f$在$p$点局部坐标系$(U;x^i)$下的像点$q$的坐标关系是
    \begin{equation*}
        y^\alpha = f^\alpha(x^1,\cdots,x^m), \quad 1 \leqslant \alpha \leqslant n .
    \end{equation*}
    因切映射$f_{*}$在点$p$是非退化的，那么可知Jacobi矩阵至少存在一个$m\times m$的
    子矩阵的行列式非零，不妨假设
    \begin{equation*}
      \left. \frac{\partial ({f^1,\cdots,f^m})}{\partial ({x^1,\cdots,x^m})} \right|_{x^i=0} \neq 0.
    \end{equation*}
    因$f$是光滑函数，此行列式至少是连续的，所以存在$p$点附近一个充分小邻域，在此小邻域范围内上式恒不为零。

    可以将$(U;x^i)$的$m$维坐标卡延拓至$n$维，即假设有$\R^{n-m}$中的一个充分小方体
    \begin{equation*}
        I:=\{(x^{m+1},\cdots,x^n) \mid \ \max(|x^{m+1}|,\cdots,|x^n|) < \delta \} ,
    \end{equation*}
    其中$\delta$是一个充分小正实数。
    我们构造映射$\tilde{f}:U\times I \to V$，其坐标表达式如下
    \eq{\label{chdm:eqn_tmp88}
        \begin{cases}
            (y\circ \tilde{f})^i = f^i(x^1,\cdots,x^m), & 1 \leqslant i \leqslant m; \\
            (y\circ \tilde{f})^\alpha =x^\alpha+ f^\alpha(x^1,\cdots,x^m), & 1+m \leqslant \alpha \leqslant n .
        \end{cases}
    }
    不失一般性，以设$p$点坐标是原点$0$；如果不是，平移一下原点即可。
    我们将方体$I$的原点和$p$结合在一起，构成$U\times I $的新原点$\tilde p:=(p,0)$。
    很明显映射$\tilde{f}$在新原点$\tilde{p}$的Jacobi行列式不为零，也就是切映射$\tilde{f}_{*}$是非退化的；
    根据定理\ref{chdm:thm_inv-in-M}可知映射$\tilde{f}$在$\tilde{p}$点充分小邻域内是微分同胚。

    现在我们舍弃$V$中原有坐标系$\{y^\alpha\}$（如果$\{y^\alpha\}$与式\eqref{chdm:eqn_tmp98}相同，
    就不用舍弃了，直接用！）；
    利用微分同胚$\tilde{f}$带来的坐标当成开集$V$的新局部坐标系$(V;z)$，可表示为
    \eq{\label{chdm:eqn_tmp98}
    \begin{cases}
        (z\circ \tilde{f})^i = x^i, & 1 \leqslant i \leqslant m; \\
        (z\circ \tilde{f})^\alpha =x^\alpha, & 1+m \leqslant \alpha \leqslant n .
    \end{cases}
    }
    在此局部坐标系内，$\tilde{f}$是$(U\times I;x)$到$(V;z)$的恒等映射。
    由映射$\tilde{f}$定义\eqref{chdm:eqn_tmp88}可知$\tilde{f}|_{U\times\{0\}}$满足定理中
    要求的映射$f|_U$；也就是，坐标系\eqref{chdm:eqn_tmp98}具体表达为
    \eq{
        \tilde{f}(x^1,\cdots,x^m,0,\cdots,0)
        = (x^1,\cdots,x^m,0,\cdots,0).
    }
    此定理说明，通过坐标变换满秩映射$f$可以将$U$单一地映入$N$的一个
    局部坐标面内。
\end{proof}


\subsection{浸入、嵌入与淹没}\label{chdm:sec_immerse-embed}
现在引入几个定义。设光滑流形$M$和$N$的维数分别是$m$和$n$；
$\phi:M\to N$是光滑映射。将$p\in M$点的Jacobi矩阵（见注解\ref{chdm:rek_jacobi}）
秩记为$\operatorname{rank}_p \phi$，可以考察关于秩的几种重要情形。
需要指出的是流形上的秩未必逐点相同，但考虑到映射是连续的，
肯定存在$p\in M$点附近的开邻域$U$，在$U$内任意点的秩都相同。

(1) 若$m=n=\operatorname{rank}_p\phi$，则$\phi$是局部微分同胚的。

(2) 若$m\leqslant n, \operatorname{rank}_p\phi=m$，也就是映射$\phi$说在$p$点某个邻域内是局部单一的，
则称$\phi$点$p$是\textbf{局部浸入}的。若
映射$\phi$在$M$上每一点都浸入，则称为\textbf{浸入映射}（immerse）。

(3) 设$\phi:M\to N$是浸入的，若$\phi$是整体单一的，则称为\textbf{嵌入映射}（embed）。

(4) 设$\phi:M\to N$是嵌入的，若$\phi$从$M$带到$\phi(M)$的拓扑结构与$\phi(M)$从$N$上得到
的诱导拓扑结构（见定义\ref{chtop:def_induced-top}）相同，
“相同”是指由$\phi$带来的拓扑和从$N$诱导得到的拓扑相互同胚；
则称之为\textbf{正则嵌入映射}（regular embedding）。

(5) 若$m> n, \operatorname{rank}_p\phi=n$，则称$\phi$点$p$是\textbf{局部淹没}的（submersion）。
%如果映射$\phi$在$M$上每一点都淹没，则称之为\textbf{淹没映射}。


浸入映射只要求局部单一，并不要求大范围单一；
因此浸入映射的像可能有自交点。
嵌入映射则要求在整个流形范围内是单一的，没有自交点。

%单纯的嵌入映射$\phi$从$M$带入到$\phi(M)$的拓扑未必与$\phi(M)$从$N$上得到
%的诱导拓扑结构相同，一般说来会细于诱导拓扑。

有了浸入和嵌入便可有了子流形定义：
\begin{definition}
    有$m$维光滑流形$M$和光滑流形$N$间的光滑映射$\phi:M\to N$；
    若$\phi$是浸入（嵌入、正则嵌入）的，则称$(\phi,M)$是$N$的一个$m$维
    \textbf{浸入（嵌入、正则嵌入）子流形}。
\end{definition}


\begin{definition}\label{chdm:def_inclusion-map}
    设有集合$A$和$B$，并且有$A\subset B$，
    映射$\imath: A\to B;\ \imath(x) = x, \ \forall x\in A. $
    那么，称映射$\imath$为\textbf{包含映射}（inclusion map）。
\end{definition}
包含映射一定是嵌入映射；但反之未必，因为嵌入映射的值域未必与定义域有交集。

\subsection{浸入与嵌入区别}
我们先叙述几个例题。
\begin{example}\label{chdm:exm_opensub}
    \textbf{开子流形}。设$U$是光滑流形$N$的一个开子集，将流形$N$的光滑结构
    限制在$U$上，变得到$U$的一个光滑结构，使之成为与$N$同维数的光滑流形。
    包含映射$(\imath,U)$是$N$的一个正则嵌入子流形。
    则称$(\imath,U)$是$N$的开子流形。
\end{example}

\begin{example}\label{chdm:exm_closesub}
    \textbf{闭子流形}。设$(\phi,M)$是$N$的一个光滑嵌入子流形，
    如果(1)$\phi(M)$是$N$的一个闭子集。 (2)$\forall p\in \phi(M)$，$N$中存在
    局部坐标系$(U;x^i)$使得$\phi(M)\cap U$是
    由方程$x^{m+1}=x^{m+2}=\cdots=x^n=0$定义的，
    那么，称$(\phi,M)$是$N$的闭子流形。
    比如，二维球面连同包含映射就是三维空间的一个闭子集。
\end{example}


下面给出两个十分相近的函数，它们都是从$\R$到$\R^2$的映射。
\begin{example}\label{chdm:exm_iec1}
    浸入，非嵌入例子。
    \eq{\label{chdm:eqn_tmp-f-immerse}
        F(t)=\left(2\cos\Big(t-\frac{\pi}{2}\Big),\ \sin2\Big(t-\frac{\pi}{2}\Big) \right) .
    }
    当$t$取值从$-\infty$到$+\infty$时，式\eqref{chdm:eqn_tmp-f-immerse}在原点（$(0,0)$）处无穷多次自相交。
    所以$(F,\R)$是$\R^2$的浸入子流形，不是嵌入子流形。
\end{example}

\begin{example}\label{chdm:exm_iec2}
    浸入，且嵌入，但非正则嵌入例子。
    \eq{\label{chdm:eqn_tmp-g-embed}
      G(t)=\left(2\cos\Big(2\arctan t+\frac{\pi}{2}\Big),\ \sin2\Big(2 \arctan t+\frac{\pi}{2}\Big) \right) .
    }
        当$t=0$时，$G(0)=(0,0)$。当$t\to \pm \infty$时，曲线无限接近原点，
    但没有到达原点，即$G(t\to \pm \infty)\to (0,0)$；没有自交点。
    所以$(G,\R)$是$\R^2$是嵌入子流形，当然它也是浸入子流形。

    由$G$带入$G(\R)$的是一条不自相交的曲线，这个拓扑是一维的，同胚于实数轴。
    我们再来看诱导拓扑，参见定义\ref{chtop:def_induced-top}。
    取二维平面$\R^2$中的通常拓扑：开圆盘。圆心在原点的开圆盘$C$，无论半径取多小
    （必须是个有限小的确定数值，不能是零，也不能取无限趋近于零的极限），首先
    会包含$-\epsilon < t< \epsilon$的$G(t)$曲线；
    其次当$t\to \pm \infty$时，$G(t\to \pm \infty)\to (0,0)$，曲线
    都会进入开圆盘$C$。很明显利用开圆盘得到的诱导拓扑不是一维的实线段，
    与映射$G$带来的拓扑不同胚，因此$G$不是正则嵌入。
\end{example}


\begin{example}\label{chdm:exm_iec3}
    浸入，且嵌入，但非正则嵌入又一例。
    定义映射$H:\R\to \R^2$如下
    \eq{
        H(t)=\begin{cases}
           \left(\frac{3}{t^2},\sin\pi t \right), & 1 \leqslant t < +\infty, \\
           \left(0,\ t+2\right), & -\infty < t \leqslant -1.
        \end{cases}
    }
    假定$H(t)$在定义区间$[-1,+1]$内是连接$(3,0)$和$(0,1)$两点的光滑曲线，且与上式给出的两个部分无交点。
    映射$(H,\R)$是整体单一浸入到$\R^2$的曲线，没有自交点；所以它是浸入的，且是嵌入的。

    带入拓扑是由$H(t)$描述的一维曲线，同胚于实数轴。
    下面分析诱导拓扑，不难发现$y$轴的$[-1,+1]$部分是曲线$H(t)$的一部分，
    记为$I_y$；显然$I_y$是诱导拓扑中的一部分，且是一维直线段。
    当$t\to +\infty$时，曲线会围绕着$x$轴上下摆动，从右侧无限接近$I_y$部分。
    与上一例相似，取$\R^2$中的通常拓扑——开圆盘$C$，
    且令开圆盘的圆心在$I_y$上。
    但是无论开圆盘半径取得多么小（必须是有限小的确定值），
    当$t\to +\infty$时，曲线$H(t\to +\infty)$都会进入开圆盘$C$；
    所以$H(t\to +\infty)$也是诱导拓扑的一部分，
    故带入拓扑与诱导拓扑不是同胚的。因此，映射$H$不是正则嵌入的。
\end{example}

我们再叙述几个定理，省略证明过程，可请参考文献\parencite[\S 1.3]{cc2001-zh}。
\begin{theorem}\label{chdm:thm_embed-1}
    设$(\phi,M)$是流形$N$的浸入子流形，若$M$是紧致的，则$(\phi,M)$是正则嵌入。
\end{theorem}

\begin{theorem}\label{chdm:thm_embed-2}
    设$(\phi,M)$是$n$光滑流形$N$的$m$维浸入子流形。
    $(\phi,M)$是正则嵌入的充分必要条件是：
    它是$N$的一个开子流形的闭子流形。
\end{theorem}

\begin{theorem}\label{chdm:thm_embed-3}
    设$(\phi,M)$是$n$光滑流形$N$的$m$维浸入子流形。
    $(\phi,M)$是正则嵌入的充分必要条件是：
    对任意点$p\in M$，在$N$中存在$q=\phi(q)$的局部坐标系$(V;y^i)$，
    其中$y^i(q)=0$；还要使得$\phi(M)\cap V$是
    由$y^{m+1}=y^{m+2}=\cdots=y^{n}=0$定义的。
\end{theorem}

对比定理\ref{chdm:thm_immerse-1}与定理\ref{chdm:thm_embed-3}，可以十分清晰地看到浸入与正则嵌入的区别。
浸入只是局部上的要求，即$\forall p\in U\subset M$，$N$中都存在$\phi(p)$点附近
的开邻域$\phi(U)\subset V\subset N$，在邻域$V$内坐标分量$y^i=0(i>m)$；对邻域$U$外的$M$中的点
如何在$N$中表示不作任何要求。
正则嵌入则是整体上的要求，即$\forall p\in U\subset M$，$N$中都存在$\phi(p)$点附近
的开邻域$\phi(U)\subset V\subset N$，在邻域$V$内坐标分量$y^i=0(i>m)$；对邻域$U$外的$M$中的点
如何在$N$中表示也作要求，即$M-U$中的点不能映入$V$中。
若只在$p$点局部邻域$U$内看问题的话，则浸入和正则嵌入没有差别，
换句话说：浸入{在局部上}必定是正则嵌入的。
例\ref{chdm:exm_iec1}、\ref{chdm:exm_iec2}、\ref{chdm:exm_iec3}具体地表明了
浸入、嵌入、正则嵌入的差别，尤其请参阅后两道例题中带入拓扑和诱导拓扑的分析。


\section{切矢量丛}\label{chdm:sec_tangent-bundles}
两个流形（笛卡尔）积是一个很基本的概念，见\S \ref{chdm:sec_cartensian-product}。
例如，定义在微分流形$M$上实函数$f(x)$的像$(x,f(x))$便是从$M$到（笛卡尔）积
流形$M\times \R$的一个映射；这个概念可以进一步推广成{纤维丛}；
本节只叙述相对简单的纤维丛——切矢量丛。


%切矢量丛是一个非常基本的概念，这里只给出结果；详尽完整论述可参见\textcite[pp. 63-106]{spivak-dif-1}或类似文献。

设$M$是$m$维光滑流形，$p\in M$点的切空间记为$T_pM$，记
\eq{\label{chdm:eqn_tangent-bundle-set}
    TM:=\bigsqcup_{p\in M}T_pM=\{(p,v):p\in M,\ v\in T_pM\} .
}
下文将 $(p,v)$ 简记为 $v$，其所在点由 $v\in T_pM$ 确定。$TM$ 是流形 $M$ 上全体切矢量的集合。我们要给它赋予一个拓扑结构和微分结构，
使它成为一个$2m$维的光滑流形。
我们将略去如下证明（可见\parencite[\S 3.1]{chenwh2001}、\parencite[\S 12.1]{tu-IM-2011}）：
$TM$是Hausdorff空间；$TM$存在开覆盖$\bigcup _{\alpha\in I} \pi^{-1} (U_\alpha)$，
并且这族开覆盖构成$TM$的可数拓扑基。

虽然省略了一些证明，但还是要建立$TM$拓扑的。
定义映射$\pi:TM\to M$如下：
\eq{
    \pi(v) = p; \quad \forall v\in T_pM,\quad p\in M.
}
这样，对每一个点$p\in M$，都有$\pi^{-1}(p)=T_pM$；注意这里的$\pi^{-1}$不是指
逆映射，是指逆像的集合；映射$\pi$不是单射，无逆；容易看出映射$\pi$是满射。

设微分流形$M$有微分结构$\mathscr{A}=\{(U_\alpha,\varphi_\alpha;x^i)|\alpha \in I \}$，令
\eq{
    \pi^{-1}(U_\alpha)=\bigsqcup_{p\in U_\alpha}T_pM .
}
当$\alpha$遍历整个指标集$I$后，
便有$\bigcup _{\alpha \in I}  \pi^{-1} (U_\alpha) = TM$。
很明显$\pi^{-1} (U_\alpha)$是$TM$的子集；
因$T_pM$是$m$维的，而点$p\in U_\alpha$（$U_\alpha$是$m$维）又是可变
动的，所以可知$\pi^{-1} (U_\alpha)$是$2m$维；进而$TM$也是$2m$维的。
这个子集中的元素自然是切矢量；
前面我们说过建立微分流形以及坐标卡的过程无法摆脱局部坐标系（见\S \ref{chdm:sec_tmpdmcon}），
我们当然选择自然坐标基矢来描述$\pi^{-1} (U_\alpha)$中的切矢量。


设$(x_p^1,\cdots,x_p^m)$是$p$点在映射$\varphi_\alpha$作用下的坐标；
矢量$v\in \pi^{-1} (U_\alpha)$可以在$p$点自然切
基矢$\left.\left(\frac{\partial }{\partial x^i}\right) \right | _p$下
展开，其分量是$(v^1,\cdots,v^m)$。
$\forall \alpha \in I$，
定义映射$\psi_\alpha: \pi^{-1} (U_\alpha) \to U_\alpha \times \R^m
\overset{\varphi_\alpha}{\longrightarrow} \varphi_\alpha(U_\alpha) \times \R^m $
（其中$\varphi_\alpha(U_\alpha) \subset \R^m$维数是$m$）。
\eq{\label{chdm:eqn_tb_atalas}
    \psi_\alpha(v):=\bigl(x^1(p),\ldots,x^m(p),v^1,\ldots,v^m\bigr),
}
其中 $p\in U_\alpha$，$v=v^i\p_i|_p\in T_pM$。
因$(x_p^1,\cdots,x_p^m)$与$(v^1,\cdots,v^m)$都是可变动的，所以上述
定义更好地说明了$\pi^{-1} (U_\alpha)$是$2m$维的。
不难看出映射$\psi_\alpha$是双射，其逆存在；
我们可以规定$\psi_\alpha$和$\psi_\alpha^{-1}$是连续映射，
相当于规定了$\pi^{-1} (U_\alpha)$中的开集；这样，
式\eqref{chdm:eqn_tb_atalas}中的映射便是拓扑同胚的，
它将$\varphi_\alpha(U_\alpha) \times \R^m $中的（通常）拓扑
携带到$\pi^{-1} (U_\alpha)$中。
很明显$(\pi^{-1} (U_\alpha),\psi_\alpha)$是$TM$一个坐标卡，
这样我们建立了$TM$的拓扑结构。
{\footnote{建立$M$的拓扑过程大体如下：(1)找到$M$的一个覆盖集合$\mathcal{C}$；
        (2)将集合$\mathcal{C}$与$\R^m$中的开子集建立双射关系$\sigma$，
        要求$\sigma$和$\sigma^{-1}$都连续；
        (3)双射$\sigma$可将$\R^m$中拓扑（即开集概念）携带到$M$；
        $\sigma$便是拓扑同胚映射。这样，$M$成为拓扑流形。

    需注意：如果$M$中原来有拓扑结构$\mathcal{T}$，这个拓扑未必与$\R^m$中开集
存在条件(2)中的$\sigma$；如果存在，那么$(M,\mathcal{T})$本身便是拓扑流形。}}


我们开始建立$TM$上的微分结构；只需验证坐标卡$(\pi^{-1} (U_\alpha),\psi_\alpha)$是$C^\infty$相容的。
实际上，$\pi^{-1} (U_\alpha) \cap \pi^{-1} (U_\beta) \neq \varnothing {\Leftrightarrow}
U_\alpha \cap U_\beta \neq \varnothing$。
当$U_\alpha \cap U_\beta \neq \varnothing$时，设分别存在
坐标卡$(U_\alpha;x^i)$和$(U_\beta;\tilde{x}^i)$；那么，它们间的坐标变换是$C^\infty$的，
即Jacobi矩阵$\frac{\partial \tilde{x}^j}{\partial {x}^i}$是$\{x^j\}$的$C^\infty$函数，
同样，其逆矩阵$\frac{\partial {x}^i}{\partial \tilde{x}^j}$也是$\{\tilde{x}^i\}$的$C^\infty$函数。
设$p \in U_\alpha \cap U_\beta$，则式\eqref{chdm:eqn_tb_atalas}中的$v$在$p$点的
坐标表达式有两个等价描述如下：
\eq{
    v = v^i \left. \left(\frac{\partial }{\partial x^i}\right) \right | _p
        = \tilde{v}^j \left. \left(\frac{\partial }{\partial \tilde{x}^j}\right)  \right | _p
}
不难得到两者分量的变换关系：
\eq{\label{chdm:eqn_tbv}
    v^i = \tilde{v}^j \left. \frac{\partial {x}^i}{\partial \tilde{x}^j}  \right | _p
    {\quad \Leftrightarrow \quad}
    \tilde{v}^j = {v}^i \left. \frac{\partial \tilde{x}^j}{\partial {x}^i}  \right | _p .
}
上述坐标变换关系中的Jacobi矩阵是$C^\infty$的，坐标$(v^1,\cdots,v^m)$本身也是$C^\infty$的，
那么坐标变换$v^i \leftrightarrow \tilde{v}^i$是$C^\infty$的；
而$U_\alpha \cap U_\beta$自身的坐标变换$x^i \leftrightarrow \tilde{x}^i$本来就是$C^\infty$的；
$TM$坐标是由两者结合得到的$(x^1,\cdots,x^m,v^1,\cdots,v^m)$，它的坐标变换自然也是$C^\infty$的。
由\pageref{chdm:def_Dmanifold}页
定义\ref{chdm:def_Dmanifold}中第二条和命题\ref{chdm:thm_biggestDiffStruc}可知，
这个坐标卡集合是{相容的}，那么便建立了$TM$的$C^\infty$微分结构，从而它是微分流形了。
$TM$便是切矢量丛，简称切丛。

更多内容请参见第\ref{chfb}章。


\section{切矢量场}\label{chdm:sec_tangent-vector-field}
在\S\ref{chdm:sec_tangentspace}中，我们已经介绍了微分流形$M$上某确定点$x$切矢量
定义\ref{chdm:def_tang-vec-inManifold}，这个定义可以推广到整个流形，那便是切矢量场。
\begin{definition}\label{chdm:def_tangent-vector-field}
    在微分流形$M$的每一点$x$都指定一个切矢量$v(x)\in T_xM$，称$v$是$M$上
    的一个\textbf{切矢量场}。
\end{definition}
%由定义不难看出切矢量场是这样一个映射$v:M\to TM$。考虑到切丛$(TM,M,\pi)$和
%矢量丛截面定义\ref{chfb:def_vf-section}，则可知$\pi \circ v = \id$；因此
%切矢量场$v$是切丛$TM$的一个截面。


设$m$维微分流形$M$有局部坐标系$(U,\varphi;x)$，很明显，
$\{(\frac{\partial }{\partial x^i})\}$（其中$1\leqslant i \leqslant m$）是
开集$U$上的切矢量场，但一般说来它们不是整个流形$M$上的切矢量场。
在\S\ref{chdm:sec_tangentspace}中，已经说明在确定的点$p$，
$\{(\frac{\partial }{\partial x^i})_p\}$是该点切空间自然坐标基矢；
在$U$上每一点都有定义的$\{(\frac{\partial }{\partial x^i})\}$（$1\leqslant i \leqslant m$）是
$U$上\textbf{自然标架场}，任一$U$上切矢量场$v(x)$都可以用此标架展开
\eq{\label{chdm:eqn_vfnf}
    v(x)|_U = v^i \left(\frac{\partial }{\partial x^i}\right) .
}


\begin{definition}
    若式\eqref{chdm:eqn_vfnf}中系数$v^i(1\leqslant i \leqslant m)$是$C^\infty(U)$的，
    则称$v$是$U$上\textbf{光滑切矢量场}。
    若$v^i$对流形$M$每个坐标卡中开覆盖是$C^\infty$的，则称$v$是$M$上\textbf{光滑切矢量场}；$M$上
    全体光滑切矢量场记为$\mathfrak{X}(M)$。
\end{definition}
用自然方式定义$\mathfrak{X}(M)$中矢量场的加法与
数乘（$\forall u,v \in \mathfrak{X}(M),\  \forall x\in M, \  \forall \lambda \in \R$），
\eq{\label{chdm:eqn_rnumtimesv}
    (u+v)(x) := u(x) + v(x), \quad
    (\lambda \cdot v)(x) := \lambda\cdot v(x).
}
这样定义后，$\mathfrak{X}(M)$构成了实数域$\R$上的矢量空间；
由于点$x\in M$是连续的、变动的，所以这个矢量空间是无穷维的（因$x$有无穷多个值）。
除了实数域上的数量乘法，还可以定义$C^\infty(M)$上函数乘法（也
称为$C^r(M)$-线性，也见定理\ref{chdm:thm_Tensor-Characterization-Lemma}；
注意与数域$\mathbb{F}$-线性的区别）：
\eq{\label{chdm:eqn_functimesv}
    (f \cdot v)(x) := f(x)\cdot v(x); \
    \forall x\in M, \ \forall f(x) \in C^\infty(M), \ \forall v \in \mathfrak{X}(M) .
}
流形$M$（或其开子集$U$）上的$C^\infty(M)$构成一个含有幺元的交换环（请读者依照定义\ref{chtop:def_ring}验证）；
在式\eqref{chdm:eqn_functimesv}定义下（加法仍由式\eqref{chdm:eqn_rnumtimesv}中第一式来定义），
$\mathfrak{X}(M)$构成了环$C^\infty(M)$上的矢量空间，也称为$C^\infty(M)$-模；
模的定义见\pageref{chtop:def_module}页\ref{chtop:def_module}。

\subsection{切矢量场几何定义}\label{chdm:sec_vdef2}
设有微分流形$M$及其局部坐标系$(U,\varphi;x)$，
再设$v\in \mathfrak{X}(M)$，及$f(x) \in C^\infty(M)$；
我们先给出切矢量场作用标量函数上的定义：
\eq{\label{chdm:eqn_vecffunc}
    \bigl(v(f)\bigr) (x)  :=  v(x) \bigl( f(x)\bigr) .
}
利用式\eqref{chdm:eqn_vfnf}，得
\eq{
    v(x) f = v^i \frac{\partial f }{\partial x^i} =
    v^i \frac{\partial f\circ \varphi^{-1} }{\partial x^i}.
}
因$v^i, f \in C^\infty(M)$，所以由上式可知光滑切矢量场作用在光滑标量场上得到
一个光滑标量场，也就是可以把光滑切矢量场看成这样一个映射$v:C^\infty(M)\to C^\infty(M)$；
上面局部坐标系是任选的，所以这种看法与局部坐标系无关。
此映射自然满足一些性质，见下面定理。


\begin{theorem}\label{chdm:thm_vf-def2}
    依式\eqref{chdm:eqn_vecffunc}定义的$v\in \mathfrak{X}(M)$是这样一个映射$v:C^\infty(M)\to C^\infty(M)$，
    并且它满足如下条件：

    (1) 线性性：$\forall g,h \in C^\infty(M) $ ， $\forall \lambda \in \R$，有
    $v (g+\lambda h) = v (g)+ \lambda v (h) $；

    (2) Leibniz律：$\forall g,h \in C^\infty(M) $，有
    $v (g\cdot h) = \bigl(v(g)\bigr)\cdot h(x)+ g(x)\cdot v (h) .$

    反之，如果有映射$\alpha:C^\infty(M)\to C^\infty(M) $满足上面两个条件，
    那么存在唯一$v\in \mathfrak{X}(M)$使得
    $v(f)=\alpha(f),\ \forall f\in C^\infty(M)$。
\end{theorem}
这个定理说明：如果映射满足上述两个条件，那么它便是一个切矢量场，从而这两个条件
可以看作是切矢量场的几何（或第二）定义；与切矢量定义\ref{chdm:def_tang-vec-inManifold}相对应；
切矢量场的第二定义不要事先给定局部坐标系，显得更几何化。

定理的前一半证明很容易，直接逐点验证即可（留给读者当练习）。
后半部分证明略显复杂，为此
我们先证明一个命题，也被称为\textbf{局部性定理}。


\begin{proposition}\label{chdm:thm_vf-local}
    设$f,c\in C^\infty(M)$，在开集$U\subset M$上有$f|_U = c|_U$。
    如果映射$v:C^\infty(M)\to C^\infty(M) $满足定理\ref{chdm:thm_vf-def2}中的两个条件，
    那么必然有$v(f)|_U = v(c)|_U$。
\end{proposition}
\begin{proof}
    利用第二个条件，令$g=1,\ h=1$，则有
    $v(1\cdot 1)=1\cdot v(1)+1\cdot v(1) {\Rightarrow}
    v(1)=2\cdot v(1){\Rightarrow} v(1)=0$，
    进而可以得到$v(const)=0$；
    这说明$v$将常数标量场映射为零标量场。

    $\forall p\in U$，选择点$p$的一个足够小邻域$V$，使得$V\subset \overline{V} \subset U$；
    令$h\in C^\infty(M)$是满足命题\ref{chdm:thm_exp1UVM}中条件的标量场，那么
    在整个流形$M$中有$(f-c)\cdot h \equiv 0$；将映射$v$作用在这个等式两端，
    并利用Leibniz律，有
    \begin{equation*}
        0=v\bigl((f-c)\cdot h\bigr)=\bigl(v(f-c)\bigr)\cdot h + (f-c)\cdot v(h)
        {\xLongrightarrow[\text{在$V$上}]{\text{限制}}}
        \bigl(v(f-c)\bigr)|_V \cdot h|_V = 0 .
    \end{equation*}
    因$h$限制在$V$上是恒为$1$的，（并利用线性性）
    所以有$v(f)|_V = v(c)|_V$。又因点$p$是任意的，所以必然有$v(f)|_U = v(c)|_U$。
\end{proof}

\begin{remark}\label{chdm:rmk_local}
定理\ref{chdm:thm_vf-def2}中的映射$\alpha$是定义在整个流形$M$上的；
一般说来，流形是不存在全局（即大范围、整体）坐标系的。
只有限制在局部开子集$U\subset M$上，才可能存在一个坐标系就能覆盖住$U$的情形，
进而使用局部坐标来描述切矢量。
定理\ref{chdm:thm_vf-def2}中的两个条件是针对整个流形$M$的，我们
不能随便将其应用到开子集$U$中，或者说切矢量在$U$上还没有定义；
我们通过命题\ref{chdm:thm_vf-local}诱导出开子集$U$上的一个“切矢量”定义，
并证明它满足“局部性”（即命题\ref{chdm:thm_vf-local}中描述的属性）。
没有局部性定理，也就没有局部坐标系相应映射的定义，
比如式\eqref{chdm:eqn_aflg}。在后面的内容中还会有类似的
局部性定理，如外微分、联络局部性定理等等。

如此叙述，是因为有些属性在整体与局部并不相同。
比如紧致性，实直线是非紧的，而其子集$[0,1]$是紧致的。
“局部性定理”都指出切矢量定义
（以及后面将会讲到的外微分、联络等，它们都可以算“导数”的推广）
在整体（流形$M$）与局部（开子集$U$）的性质是相同的。
\end{remark}

$\forall p\in U$，选择点$p$的一个足够小邻域$V$，使得$V\subset \overline{V} \subset U$；
取$f\in C^\infty(U)$，根据命题\ref{chdm:thm_Flocal-equiv-Fglobal}可以将$f$延拓至
整个流形$M$上的$\hat{f}\in C^\infty(M)$，并且有$f|_V= \hat{f}|_V$。
根据局部性定理，我们可以将$\alpha(f)$在$p\in U$点的局部值定义为
\eq{\label{chdm:eqn_aflg}
    \bigl(\alpha(f)\bigr)(p) := \bigl(\alpha(\hat{f})\bigr)(p) .
}
定义号“$:=$”右端是$C^\infty(M)$中的标量场，本来就有意义。
需要注意的是，这个定义与开集$V$选取无关；
可以将$V$换成任意其它开集$V'$，只要满足$V' \subset \overline{V'} \subset U$即可；
证明过程与上面完全相同。
也和函数$\hat{f}$的选取无关，可以将$\hat{f}$换成任意$\tilde{f}$，
只需满足$f|_V=\hat{f}|_V=\tilde{f}|_V$即可（有无穷多种$\tilde{f}$满足这个关系）。
在通常的微积分中，我们谈及函数$f(x)$在$p$点可导属性时，只需涉及$p$点一个小邻域即可，
无需知道整个实数轴上$f(x)$的性状；上面局部性定理告诉我们这一性质在微分流形中
也成立（此性质同样适用于后面的外微分、联络等）。

容易验证由式\eqref{chdm:eqn_aflg}诱导出的局部映射$\alpha:C^\infty(U)\to C^\infty(U)$同样
满足定理\ref{chdm:thm_vf-def2}中的两个条件。下面开始证明定理\ref{chdm:thm_vf-def2}。

设流形$M$有相容性坐标覆盖$(U_i,\varphi_i;x^j)$，对于任意点$p\in M$，它必然属于
某个开覆盖$U_i$；在此点，我们定义一个映射$v:C^\infty_p \to \R$如下
\eq{\label{chdm:eqn_vfgl}
    \bigl(v(f)\bigr)(p) = \bigl(v(p)\bigr)(f) :=
      \bigl(\alpha({f})\bigr)(p) =\bigl(\alpha(\hat{f})\bigr)(p) , \quad \forall f\in C^\infty_p .
}
由于局部映射$\alpha:C^\infty(U)\to C^\infty(U)$同样
满足定理\ref{chdm:thm_vf-def2}中的两个条件；不难证明
在定义\eqref{chdm:eqn_vfgl}的意义下，这两个条件全部限制在
点$p$处取值，即（证明留给读者，极易。
其中$\forall g,h \in C^\infty_p,\ \lambda \in \R $）
\eq{\label{chdm:eqn_tmpvfg}
\begin{aligned}
    \bigl(v(p)\bigr) (g+\lambda h) &= \bigl(v(p)\bigr) (g) + \lambda\cdot \bigl(v(p)\bigr) (h) , \\
    \bigl(v(p)\bigr) (g\cdot h) &= \bigl(v(p)(g)\bigr) \cdot h + g\cdot\bigl(v(p)\bigr) ( h) .
\end{aligned}
}
上两式说明，由式\eqref{chdm:eqn_vfgl}定义的$v(p)$符合切矢量定义\ref{chdm:def_tang-vec-inManifold}，
因此$v(p)$是流形$M$在点$p$处的一个切矢量。由于点$p$是任意选取的，于是$M$中每一点都可以定义类似
切矢量，这样构造的$v$是$M$上的{切矢量场}。

仍需证明$v$的光滑性。取局部坐标系$(U_i,\varphi_i;x^j)$，将$v$展开

\eq{\label{chdm:eqn_tmpadf23}
    v(p) = v^j(p) \left.\left(\frac{\partial }{\partial x^j}\right) \right|_p
      = \bigl(v(p)\bigr)(x^j) \left.\left(\frac{\partial }{\partial x^j}\right) \right|_p
      = \bigl(\alpha(p)\bigr)(x^j) \left.\left(\frac{\partial }{\partial x^j}\right) \right|_p
}
坐标函数$x^j$是光滑标量场，故而$\alpha(x^j)$也是光滑标量场，由定义\eqref{chdm:eqn_aflg}可知
其限制在$p$点取值后，$v^j(p)$也是$C^\infty_p$的；又因$p$是任取的，所以$v^j$是$C^\infty(M)$的。
这便说明了$v$在$M$上的光滑性。
切矢量场光滑性的证明不依赖于某特定局部坐标系。定理\ref{chdm:thm_vf-def2}证毕。

\begin{remark}
    定理\ref{chdm:thm_vf-def2}的主要目的是提供一个（表面上）不用局部坐标系的方式
    定义切矢量场，也就是由此定理能导出定义\ref{chdm:def_tangent-vector-field}；
    从而说明两个定义是等价的。

    定理\ref{chdm:thm_vf-def2}中两个条件是针对整个流形而言的，不能随便用到局部坐标系。
    选取局部坐标系后，通过局部性定理以及定义式\eqref{chdm:eqn_vfgl}，把原来
    只对整体流形$M$成立的定义诱导出了局部坐标系相应定义；而且容易证明
    此定义满足式\eqref{chdm:eqn_tmpvfg}，即诱导出的定义符合切矢量\ref{chdm:def_tangent-vector-field}的要求。
    这样，便从第二定义导出了第一定义。
\end{remark}


\subsection{Poisson括号}\label{chdm:sec_poisson-lie}
设$u,v \in \mathfrak{X}(M)$，$\forall f\in C^\infty(M)$，有
\eq{\label{chdm:def_poisson-bracket}
    \left[u,v\right] (f) := u\bigl(v(f)\bigr) - v\bigl(u(f)\bigr) .
}
由上式定义的Poisson括号（也称为对易子）是流形$M$上一个光滑矢量场，即$[u,v] \in \mathfrak{X}(M)$。
证明并不困难，首先从定义式易看出：$\left[u,v\right]:C^\infty(M)\to C^\infty(M)$；
剩下只需验证它满足定理\ref{chdm:thm_vf-def2}中的两个条件，
我们先验证线性性（其中$\forall f,g \in C^\infty(M),\ \lambda \in \R $）：

\begin{align*}
    &\left[u,v\right] (f + \lambda g) = u\bigl(v(f + \lambda g)\bigr) - v\bigl(u(f + \lambda g)\bigr)
    =u\bigl(v(f)+\lambda v(g)\bigr)- v\bigl(u(f)+\lambda u(g)\bigr) \\
    &=u\bigl(v(f)\bigr)+ \lambda u\bigl(v(g)\bigr) - v\bigl(u(f)\bigr) - \lambda v\bigl( u(g)\bigr)
    =\left[u,v\right] (f ) + \lambda \left[u,v\right] ( g) .
\end{align*}
Leibniz律的验证略显麻烦，我们先计算下式：
\begin{align*}
    u\bigl(v(f \cdot g)\bigr) &= u\bigl( v(f) \cdot g+ f \cdot v( g)\bigr)
      = u\bigl( v(f) \cdot g\bigr) + u\bigl( f \cdot v( g)\bigr) \\
      &= u\bigl( v(f) \bigr) \cdot g + v(f) \cdot u\bigl(  g\bigr)
      + u\bigl( f \bigr) \cdot v(g) + f \cdot u\bigl( v( g)\bigr) .
\end{align*}
将上式中$u$和$v$互换，并作差，得
\begin{align*}
    \left[u,v\right] (f \cdot g) &= u\bigl( v(f) \bigr) \cdot g + f \cdot u\bigl( v( g)\bigr)
      -v\bigl( u(f) \bigr) \cdot g - f \cdot v\bigl( u( g)\bigr) \\
      & =\left\{ \left[u,v\right](f) \right\} \cdot g + f \cdot \left[u,v\right]  (g)
\end{align*}
这便验证了Leibniz律。因此证明了Poisson括号是一个光滑切矢量场。

设$M$有局部坐标系$(U;x^i)$，可在局部坐标系将Poisson括号展开，有
\eq{\label{chdm:eqn_poisson-bracket-expand}
\begin{aligned}
    \left[u,v\right] (f) &= u^i\frac{\partial }{\partial x^i}\bigl(v^j\frac{\partial }{\partial x^j}(f)\bigr)
       - v^j\frac{\partial }{\partial x^j}\bigl(u^i\frac{\partial }{\partial x^i}(f)\bigr) \\
     &=  u^i\frac{\partial v^j}{\partial x^i} \frac{\partial f}{\partial x^j}
       + u^iv^j \frac{\partial^2 f}{\partial x^i x^j}
       - v^j\frac{\partial u^i}{\partial x^j}\frac{\partial f}{\partial x^i}
       - u^i v^j\frac{\partial^2 f}{\partial x^j x^i} \\
     &=  \left( u^i\frac{\partial v^j}{\partial x^i}
       - v^i\frac{\partial u^j}{\partial x^i} \right)  \frac{\partial }{\partial x^j} (f) .
\end{aligned}
}

\begin{theorem}\label{chdm:thm_poisson-Lie-bracket}
    光滑流形$M$上的切矢量场$\mathfrak{X}(M)$中的Poisson括号满足如下规则：

    (1) 反交换律：$\left[u,v\right] = -\left[v,u\right]$；

    (2) $\R$-线性：$\left[k u+ v,w\right] = k \left[u,w\right] +\left[v,w\right]$；

    (3) Jacobi恒等式：$\bigl[u, [v,w]\bigr]+\bigl[w, [u,v]\bigr]+\bigl[v, [w,u]\bigr]=0$。

    其中$\forall u,v,w \in \mathfrak{X}(M)$，$ \forall k \in \R$。
\end{theorem}
\begin{proof}
    前两条性质证明极易；下面给出第三条性质的证明，
    $\forall f\in C^\infty (M)$，有
    \begin{align*}
        \bigl[u, [v,w]\bigr](f) & =[u, vw-wv](f)=(uvw-uwv-vwu+wvu)(f) \\
        \bigl[w, [u,v]\bigr](f) & =(wuv-wvu-uvw+vuw)(f) \\
        \bigl[v, [w,u]\bigr](f) & =(vwu-vuw-wuv+uwv)(f)
    \end{align*}
    上面三式相加便是定理中的第三条。
\end{proof}
满足定理\ref{chdm:thm_poisson-Lie-bracket}中三个条件的运算，可以看成某种乘法，
光滑切矢量场$\mathfrak{X}(M)$便是实数域上关于Poisson括号乘积的\textbf{李代数}；
更详细内容请见例\ref{chlg:exm_LAXM}。 %\S \ref{chlg:sec_definition}。

Poisson 括号具有 $\R$-线性，但一般不具有 $C^\infty(M)$-线性。任给 $f,g\in C^\infty(M)$ 及 $u,v\in\mathfrak X(M)$，有
\eq{\label{chdm:eqn_fg-PoissonBracket}
    [f u, g v] = fg[u,v]+f\cdot u(g)\cdot v -g\cdot v(f)\cdot u .
}
计算过程并不复杂，留给读者当作练习。


\subsection{场上的诱导映射}\label{chdm:sec_induced-map-on-fields}

在\S\ref{chdm:sec_induced-map}中，引入确定点$p$的切映射、余切映射等概念；
其中有一些可以直接推广到“场”上，有一些则不可以。

设有两个微分流形$M$和$N$，维数分别为$m$和$n$（两者未必相等），存在光滑映射$\phi:M\to N$；
$M$和$N$的局部坐标分别为$(U;x^i)$和$(V;y^\alpha)$。


其中标量函数的拉回映射\ref{chdm:def_pullback-scalar}可以直接推广到“场”与“场”间的拉回：
\begin{definition}\label{chdm:def_pullback-scalar-onfield}
    从$C^\infty(N)$到$C^\infty(M)$
    的\textbf{拉回映射} $\phi^{*}$定义为
    \eq{
        \phi^{*}(g) := g \circ \phi,
        \quad \forall g\in C^\infty(N) .
    }
\end{definition}
由于标量函数的拉回映射本质上就是复合函数，对$\phi$没有太多要求；例如，
举一个极端的例子，假设两个流形分别都是实直线$\R$，
$\phi$是个常数映射 $\phi(x)=0,\ \forall x \in \R$，
这个映射显然不是单射，也不是满射。那么$\phi$定义的拉回是良定义吗？
检查拉回的定义就能看出一点问题都没有，$\phi$把（目标流形$N$）$\R$的
任何光滑函数$g$拉回为（源流形$M$）$\R$的常数函数，其值为 $g(0)$；
$\phi^*(g)$对任意$x\in M$都有定义。

仿照切矢量场定义，给出余切矢量场定义：
\begin{definition}\label{chdm:def_cotangent-vector-field}
    在微分流形$M$的每一点$x$都指定一个余切矢量$\omega(x)\in T^{*}_xM$，称$\omega$是$M$上
    的一个\textbf{余切矢量场}。一般记为$\mathfrak{X}^{*}(M)$或者$A^1(M)$。
\end{definition}

余切矢量的拉回同样可以拓展到“场”：
\begin{definition}\label{chdm:def_pullback-1form-onfield}
    余切矢量场的\textbf{拉回}逐点定义为
    \eq{
        (\phi^*\omega)_x(v):=\omega_{\phi(x)}(\phi_*v),\quad
        \omega\in\mathfrak X^*(N),\quad v\in T_xM.
    }
\end{definition}
可以看出这个定义中的点$x$是可以遍历到流形$M$中任意一点的；
$(\phi_{*}v)$只需逐点定义即可。
映射$\phi$是否单一（是否满射）并不影响$(\phi^{*}\omega)$是$M$上的“场”这一结论。
不难看出定义\ref{chdm:def_pullback-1form-onfield}很容易拓展到任意$(0,q)$型张量场。

切矢量的推前映射与以上两个定义不同，一般说来，它不能推前成“场”。
比如$\phi$不是单一的，可能有$x_1,x_2\in M$使得$\phi(x_1)=\phi(x_2)=y\in N$，
但多数情况下$\phi_*(v_{x_1}) \neq \phi_*(v_{x_2})$；这样在$y\in N$便有两个
不同的切矢量，出现多值性。
再比如$\phi$不是满射，那么$\phi_*(v_{x})$不能遍历到$N$上每个点，因此不能称为“场”。


再研究切矢场 $v$。确定一个 $v$ 只需对其作用于任意 $f$ 的结果做出要求。将 $v$ 推前即 $\phi_* v(f):=v(\phi^* f)$。这里未出现 $\phi^{-1}$，可见对一般光滑映射，只能借助切矢场给出：
\begin{definition}
    $\phi_*:T_p M\to T_{\phi(p)}N$ 称为，使 $\phi_* v:=v\circ \phi^*$。
\end{definition}

若 $\phi$ 是微分同胚，则切矢场的拉回为 $\phi^*A:=(\phi^{-1})_*A$，余切矢场的推前为 $\phi_*\omega:=(\phi^{-1})^*\omega$。一般张量场在每个切矢因子和余切矢因子上分别作相应变换，具体定义见\ref{chdm:def_push-pull-tensor}。

如果规定$\phi$是微分同胚，那么可以把推前切矢量变成“场”。
\begin{definition}\label{chdm:def_pushforward-vector-onfield}
    设 $\phi:M\to N$ 为微分同胚。切矢量场 $v$ 的\textbf{推前}由 $(\phi_*v)_{\phi(x)}:=\phi_*(v_x)$ 定义，等价地，
    \eq{\label{chdm:eqn_pushforward-vector-onfield}
        (\phi_*v)(g):=\bigl(v(g\circ\phi)\bigr)\circ\phi^{-1},\quad
        g\in C^\infty(N).
    }
\end{definition}
很明显，切映射$\phi_{*}$是流形$M$切空间$T_xM$和流形$N$切空间$T_{\phi(x)}N$间的线性映射。

\begin{proposition}\label{chdm:thm_fxxf}
    设$\phi:M\to N$是微分同胚，那么$\forall v \in\mathfrak{X}(M)$、
    $\forall \omega \in\mathfrak{X}^{*}(M)$和$f\in C^\infty(M)$，
    它们在微分同胚映射$\phi$下的变换关系是：
    $    \phi_{*} (f v) = \bigl(f\circ\phi^{-1}\bigr) \bigl( \phi_{*} v \bigr), \
        \phi^{-1*} \bigl(v(f)\bigr) = \bigl( \phi_{*} v \bigr)
        \bigl(f\circ\phi^{-1}\bigr), \
        \phi^{-1*}(f \omega) = \bigl(f\circ\phi^{-1}\bigr)\bigl(\phi^{-1*}\omega\bigr)
    $。
\end{proposition}
\begin{proof}
    由定义\ref{chdm:def_pushforward-vector-onfield}可知，$\forall g\in C^\infty(N)$，有
    \eq{
        (\phi_*(fv))(g)=\bigl(f\,v(g\circ\phi)\bigr)\circ\phi^{-1}
        =(f\circ \phi^{-1})\cdot \bigl(\phi_{*} v(g)\bigr).
    }
    需注意$f$的定义域是$M$，故需要$\phi^{-1}$操作。与之类似，可证明第二式
    \eq{
        \phi^{-1*} \bigl(v(f)\bigr) = \bigl(v(f)\bigr) \circ \phi^{-1}
        =\bigl( \phi_{*} v \bigr)  \bigl(f\circ\phi^{-1}\bigr) .
    }
    第三式证明留给读者。
    总的来说，切矢量场要用$\phi_{*}$推前；标量函数场、余切矢量场要用$\phi^{-1}$拉回。
\end{proof}


\subsection{$\phi$-相关矢量场}\label{chdm:sec_related}
定义\ref{chdm:def_pushforward-vector-onfield}将映射限定为微分同胚，其实这还能弱化。
设有两个光滑流形$M$、$N$，它们的维数分别是$m$、$n$，两者可以不等（一般$m\leqslant n$）；
再设$\phi:M\to N$是光滑映射。 %\cite[\S 1.3]{helgason-2001}
\begin{definition}\label{chdm:def_related}
    $X\in \mathfrak{X}(M)$，$Y\in \mathfrak{X}(N)$；
    若$\forall p\in M$有$\phi_{*} X_p = Y_{\phi(p)}$成立（或有
    等价式$(\phi_{*}X)(f)= (Yf)\circ\phi$，$\forall f\in C^\infty(N)$），则称
    $X$与$Y$是$\phi$-\textbf{相关矢量场}。
\end{definition}
举个简单的例子，比如流形$M$是一条直线，$N$是二维平面；$\phi$是恒等映射，
它将$M$映射为$N$中的一条直线（比如$x$坐标轴）。
$M$上（直线）自然可以定义一个矢量场$X$，映射$\phi$只是将$M$映射为$N$中的$x$轴；
在$N$中$x$轴之外的地方无法与$M$建立联系，故$\phi$不能把$X$推前成$N$中的“场”。
如果把$N$中矢量场$Y$的取值范围限定在$x$轴上（即$\phi(M)$上），
那么$X$和$Y_{\phi(p)}$间是有双射关系的，这便是
定义\ref{chdm:def_related}中描述的$\phi$-相关。

%证明摘自\parencite{warner-1983-FDMLG}定理1.55 。

\begin{theorem}\label{chdm:thm_push-Poisson-related}
    设$X_i$和$Y_i$是$\phi$-相关的（$i=1,2$），则$[X_1,X_2]$
    与$[Y_1,Y_2]$是$\phi$-相关的，
    即$\phi_{*} [X_1,X_2] = [Y_1,Y_2]\circ \phi =\left[\phi_{*}X_1,\phi_{*}X_2\right]$。
\end{theorem}
\begin{proof}
    由定义\ref{chdm:def_related}可知：
    \begin{equation*}
        Y_1\bigl( Y_2 (f)\bigr)\circ \phi = (\phi_* X_1 )\bigl( Y_2 (f)\bigr)
        =X_1\bigl( Y_2 (f\circ \phi)\bigr)  =X_1\bigl( \phi_{*} X_2 (f)\bigr) = X_1\bigl( X_2 (f\circ \phi)\bigr) .
    \end{equation*}
    由此式自然可证$\phi_{*} [X_1,X_2] = [Y_1,Y_2]\circ \phi =\left[\phi_{*}X_1,\phi_{*}X_2\right]$。
\end{proof}

由上述定理立刻可得如下定理：
\begin{theorem}\label{chdm:thm_push-Poisson}
    设$\phi:M\to N$是微分同胚，则$\phi_{*}$保持Poisson括号不变：
    \eq{\label{chdm:eqn_push-Poisson}
        \phi_{*} [X,Y] = \left[\phi_{*}X,\phi_{*}Y\right], \quad
        \forall X, Y \in \mathfrak{X}(M) .
    }
\end{theorem}
%\begin{proof}
%    设$\forall g\in C^\infty(N)$，$\forall p\in M$，记$q=\phi(p)$，那么
%%    \begin{equation*}
%%        (\phi_{*} u) (g)|_q =  (\phi_{*p} u_p) (g) = u_p\bigl(g\circ \phi(p)\bigr)
%%        = \bigl(u (g\circ \phi )\bigr)|_p .
%%    \end{equation*}
%%    因点$p$是任选的，所以对整个场而言也是正确的。
%%    利用上式可得
%    \begin{equation*}
%         \phi_{*}u  \bigl(\phi_{*}v (g) \bigr)
%         = \phi_{*}u  \bigl(v (g\circ \phi )\bigr)
%         = u  \Bigl(\bigl(v (g\circ \phi )\bigr)\circ  \phi \Bigr)
%         = u  \bigl(v (g\circ \phi ) \bigr)
%    \end{equation*}
%    利用此式可得
%    \begin{equation*}
%        \left[\phi_{*}u,\ \phi_{*}v \right] (g) =
%        \bigl(uv -vu  \bigr)(g\circ \phi )
%        =\bigl(\phi_{*}\left[u,\ v \right]\bigr) (g)
%    \end{equation*}
%    所以微分同胚的切映射保Poisson括号不变。
%    证明过程请参考式\eqref{chdm:eqn_f-cdot-g-circo-gamma}。
%\end{proof}


\begin{exercise}
	证明式\eqref{chdm:eqn_fg-PoissonBracket}。
\end{exercise}

%%%%%%%%%%%%%%%%%%%%%%%%
\section{张量场}\label{chdm:sec_tensor-fields}
切矢量场的概念自然可以推广到张量上来，见下：
\begin{definition}\label{chdm:def_tensor-field}
    在微分流形$M$每一点$x$指定一个$(p,q)$型张量
    $T(x)\in \mathcal{T}^{p}_{q} M (x)$，
    则称$T(x)$是$M$上的一个\textbf{张量场}。
\end{definition}
设点$x\in M$有局部坐标系$(U;x^i)$，张量$T$表达式为：
\eq{\label{chdm:eqn_tensor-component111}
    T=T^{i_1\cdots i_p}{}_{j_1\cdots j_q}
    \p_{i_1}\otimes\cdots\otimes\p_{i_p}\otimes
    \d x^{j_1}\otimes\cdots\otimes\d x^{j_q}.
}
\begin{definition}
    若式\eqref{chdm:eqn_tensor-component111}中系数$T^{i_1 \cdots i_p}_{\phantom{i_1 \cdots i_p}j_1\cdots j_q}$
    是$C^\infty(U)$的，
    则称$T$是$U$上\textbf{光滑张量场}。
    $U$上全体光滑切张量场记为$\mathfrak{T}^p_q(U)$；$M$上全体光滑切张量场记为$\mathfrak{T}^p_q(M)$。
\end{definition}
可以仿照切矢量场定义加法和数量乘法（见式\eqref{chdm:eqn_rnumtimesv}）使$\mathfrak{T}^p_q(M)$成为线性空间；
仿照式\eqref{chdm:eqn_functimesv}定义$C^\infty(M)$-乘法（注意，此乘法不同于数量乘法）。除此之外，还可逐点定义
张量积运算和缩并运算（见\S \ref{chmla:sec_tensor}）。如果记
\eq{
    \mathfrak{T}(M) = \bigoplus_{p,q \geqslant 0} \mathfrak{T}^p_q(M),
}
则$\mathfrak{T}(M)$成为$C^\infty(M)$模。其中，
$\mathfrak{T}^{0}_{0}(M)$便是$C^\infty(M)$；
$\mathfrak{T}^{1}_{0}(M)$便是$\mathfrak{X}(M)$；
$\mathfrak{T}^{0}_{1}(M)$是全体余切矢量场集合，也称为\textbf{一次微分式}，常用记号是${A}^{1}(M)$或$\mathfrak{X}^*(M)$；
$\mathfrak{T}^{0}_{2}(M)$中有一类特别子集称作\textbf{度量}或\textbf{度规}。

张量场也存在与定理\ref{chdm:thm_vf-def2}类似的定理，
因此张量场也可有与切矢量场类似的第二定义；两者的第二定义不需要局部坐标系，显得更几何化。
\begin{theorem}[张量场特征引理]\label{chdm:thm_Tensor-Characterization-Lemma}

    光滑流形$M$上一个多重线性映射$\mathscr{T}$如下
    \begin{equation*}
        \mathscr{T}: \underbrace{\mathfrak{X}^{*}(M) \times \cdots \times \mathfrak{X}^{*}(M)}_{p\text{个}} \times
        \underbrace{\mathfrak{X}(M) \times \cdots \times \mathfrak{X}(M)}_{q\text{个}}
        \to C^\infty(M)
    \end{equation*}
    是一个$(p,q)$型光滑张量场，等价于，映射$\mathscr{T}$对每个位置都是$C^\infty(M)$-线性的（包含$\R$-线性），即
    \eq{\label{chdm:eqn_Tensor-Characterization-Lemma}
\begin{aligned}
    &\mathscr T(\alpha_1,\ldots,\alpha_p;v_1,\ldots,fv_r,\ldots,v_q)\\
    &\quad=f\,\mathscr T(\alpha_1,\ldots,\alpha_p;v_1,\ldots,v_q)\\
    &\quad=\mathscr T(\alpha_1,\ldots,f\alpha_s,\ldots,\alpha_p;v_1,\ldots,v_q).
\end{aligned}
}
    其中，$\forall f\in C^\infty(M)$，$\forall \alpha_j\in \mathfrak{X}^{*}(M)$，$\forall v_j\in \mathfrak{X}(M)$，
    任意的$1\leqslant r \leqslant q$和$1\leqslant s \leqslant p$。
\end{theorem}
\begin{proof}
    可参考文献\parencite[\S 3.5]{chenwh2001}定理5.1，或\parencite[\S 12.3]{lee-2013-ism}引理12.24；
    或其它类似文献。
\end{proof}


很明显，这个定理有许多变种，比如（利用上面定理很容易证明如下推论）
\begin{corollary}\label{chdm:thm_Tensor-Characterization-Lemma-1}
    光滑流形$M$上一个多重线性映射$\mathscr{T}$如下
    \begin{equation*}
        \mathscr{T}: \underbrace{\mathfrak{X}^{*}(M) \times \cdots \times \mathfrak{X}^{*}(M)}_{p\text{个}} \times
        \underbrace{\mathfrak{X}(M) \times \cdots \times \mathfrak{X}(M)}_{q\text{个}}
        \to {\mathfrak{X}(M)}
    \end{equation*}
    是一个$(p+1,q)$型光滑张量场，充分必要条件是，映射$\mathscr{T}$对每个位置都是$C^\infty(M)$线性的，
    即满足式\eqref{chdm:eqn_Tensor-Characterization-Lemma}。
\end{corollary}

推前和拉回映射可拓展到$(p,q)$型张量场上。
\begin{definition}\label{chdm:def_push-pull-tensor}
    设 $\phi:M\to N$ 为微分同胚，$T$ 为 $M$ 上的 $(p,q)$ 型张量场。其\textbf{推前}由下式定义，其中 $y=\phi(x)$：
    \[
    \begin{aligned}
        &(\phi_*T)_y(\omega_1,\ldots,\omega_p;u_1,\ldots,u_q)\\
        &\quad:=T_x\bigl(\phi^*\omega_1,\ldots,\phi^*\omega_p;
        (\phi^{-1})_*u_1,\ldots,(\phi^{-1})_*u_q\bigr),
    \end{aligned}
    \]
    其中 $\omega_i\in T_y^*N$、$u_j\in T_yN$。对 $N$ 上的张量场 $S$，其\textbf{拉回}定义为 $\phi^*S:=(\phi^{-1})_*S$。在切矢和余切矢上，这些定义与前文一致。
\end{definition}

\begin{proposition}\label{chdm:thm_pptt}
    设有$N$上张量场$T$、$S$，由定义\ref{chdm:def_push-pull-tensor}可
    知$\phi^{*}(T\otimes S) = \phi^{*}(T) \otimes \phi^{*}(S)$。
\end{proposition}


\section{单参数变换群}\label{chdm:sec_One-Parameter-Transformations-Groups}
单参数变换群理论是微分流形中一个重要课题，贯穿有两个中心思想及其衍生结果，
即矢量场直化定理和单参数线性变换群理论。

我们需要用常微分方程理论来描述流形上的矢量场。为此，在本节假设常微分方程解存在且唯一，
也就是满足Lipschitz条件\cite[\S 31]{arnold-2001-ode}。


\begin{definition}
    给定$m$维光滑流形$M$上的光滑切矢量场$v$且处处非零。如果流形$M$上的光滑曲线$\gamma(t)$的
    切矢量$(\frac{\partial}{\partial t}) = v$（见式\eqref{chdm:eqn_DiffCurve-2}），
    则称曲线$\gamma(t)$是切矢量场$v$的\textbf{积分曲线}（或\textbf{轨道}）。
\end{definition}
因此，在曲线上每一点，曲线$\gamma(t)$的切矢量都由$v$规定。

设流形$M$有局部坐标系$(U;x^i)$，借用式\eqref{chdm:eqn_v-expand-on-natural-bases}，
将切矢量场$v$在局部坐标系展开，其中$v^i$是其分量。
再参考例题\ref{chdm:exm_DiffCurve}，可用分量语言来详细描述一下上述定义
\eq{\label{chdm:eqn_tmp-va2xit}
    \frac{\d x^i\circ\gamma(t)}{\d t}=v^i(x^1\circ\gamma,\cdots,x^m\circ\gamma),
    \quad \left. x^i(t) \right|_{t=t_0} = x^i_p  .
}
我们在上式中同时指定了初始条件：当$t=t_0$时，坐标$x^i(t_0)$取$p\in M$点坐标$x^i_p$。
上式是积分曲线的分量表达式，需要注意矢量场$v$是不显含参量$t$的。
根据常微分方程解存在唯一性定理\cite[\S 31]{arnold-2001-ode}可知
给定初始条件的微分系统\eqref{chdm:eqn_tmp-va2xit}解是唯一确定的；
也就是说给定矢量场$v$后，过$p$点的积分曲线是（局部）唯一的。
由于积分曲线未必在整个实数轴上有定义，所以积分曲线未必整体存在、唯一。

\begin{example}
    设有一维光滑流形：实直线$\R$，局部坐标取为$\{x\}$。其上有光滑矢量场$k x$（只有一个分量），
    其中$k$是实常数；求其过$p$点的积分曲线。将上述条件带入式\eqref{chdm:eqn_tmp-va2xit}，有
    \eq{
        \frac{\d x(t)}{\d t}= k x
        \quad {\Rightarrow} \quad
        x(t)= p e^{k(t-t_0)}.
    }
\end{example}
上面例题应该是最简单的情形之一，读者可在平面上画出积分曲线（令$p$变动）。
下面给出一个略微复杂一点的例子。
\begin{example}
    设有二维光滑流形：二维实平面$\R^2$，局部坐标取为$\{x,y\}$。
    其上有光滑矢量场$(x,ky)$（有两个分量），其中$k$是实常数；
    求其过$(p,y_0)$点的积分曲线。将上述条件带入式\eqref{chdm:eqn_tmp-va2xit}，有
    \eq{
        \begin{cases}
            \frac{\d x(t)}{\d t}& =  x \\
            \frac{\d y(t)}{\d t}&= k y
        \end{cases}
        \quad {\Rightarrow} \quad
        \begin{cases}
            x(t)& =  p e^{t-t_0} \\
            y(t)& =  y_0 e^{k(t-t_0)}
        \end{cases}
    }
积分曲线比上一例题复杂多了，需要对$k$取值范围进行分类，然后再画；
有兴趣的读者请仔细画出。
\end{example}

积分曲线不仅与出发点（即初始条件那个点）有关，还与曲线参数化有关。
\begin{example}
    设$v$是光滑流形$M$上处处非零的光滑切矢量场。
    假设$\gamma(t)$是矢量场$v$的积分曲线，
    那么$\gamma(u)$（其中$u=t+c$，$c$是实常数）
    也是矢量场$v$的积分曲线。
直接计算即可证明，
    \eq{
        \frac{\d x^i\circ\gamma(t+c)}{\d t} =
        \frac{\d x^i\circ\gamma(t+c)}{\d (t+c)}=
        \frac{\d x^i\circ\gamma(u)}{\d u}=v^i(u)=v^i(t+c).
    }
\end{example}

处处非零矢量场$v$过$p\in M$点积分曲线$\gamma(t),\ b<t<c$在$p$点小邻域内是唯一的，其中$\gamma(0)=p$。
我们将$\gamma(t)$的定义域$(b,c)$向整个实数轴延拓，使得所有符合方程\eqref{chdm:eqn_tmp-va2xit}的解
都包含在内，则称延拓后的积分曲线$\gamma(t)$为矢量场$v$的、通过点$p$的\textbf{最大积分曲线}。
需要指出的是最大积分曲线的定义域未必是整个实数轴，可能存在某些无法包含在内的点或区间（比如奇点）。

\begin{definition}
    设$v \in \mathfrak{X}(M)$，若在$p\in M$有$v(p)=0$，
    则称$p$是$v$的\textbf{奇点}。
\end{definition}
我们不研究奇点。在非奇点附近，切矢量场性状是简单的。


\begin{definition}\label{chdm:def_vector-complete}
    设$v \in \mathfrak{X}(M)$且处处非零，如果$v$的最大积分曲线$\gamma(t)$中的
    参数$t$可以取遍整个实数轴$\R$，那么称$v$是\textbf{完备}切矢量场。
\end{definition}

设矢量场$v$是完备的，则它的最大积分曲线定义域是整个实数轴。如果我们人为地挖掉某点（比如$t=1$），
那么积分曲线定义域不再是整个实数轴，矢量场也不再完备。这种点可以通过延拓来补齐，从而
使得矢量场再次完备。不能通过延拓而补上的点具有奇异性。

\begin{proposition}\label{chdm:thm_st}
    设有光滑流形$M$，处处非零的光滑矢量场$v\in \mathfrak{X}(M)$的最大积分曲线$\gamma_p(t)$过点$p \in M$，
    并且当$t=0$时，$\gamma_p(0)=p$；令$q=\gamma_p(s)$。那么有$\gamma_p(s+t)=\gamma_q(t)$。
\end{proposition}
\begin{proof}
    假设命题中所有参数取值在积分曲线$\gamma_p$定义域内。

    由命题中叙述可知，点$q\in M$也在积分曲线$\gamma_p$上，
    且其参数取$s$，即$q=\gamma_p(s)$。

    曲线$\gamma_q$是指通过$q$点的矢量场$v$的积分曲线（且$\gamma_q(0)=q$）。
    由于$q$在曲线$\gamma_p$上，并且矢量场$v$的最大积分曲线是唯一的
    （可从常微分方程解的唯一性得到）；
    所以$\gamma_q$与$\gamma_p$是同一条曲线的两种不同表述方式，差别仅在于初始点的不同。
    依命题中叙述不难得到结论：$\gamma_p(s+t)=\gamma_q(t)$。
\end{proof}

设有$m$维光滑流形$M$，对于任意点$p \in M$，
由处处非零的完备切矢量场$v\in \mathfrak{X}(M)$的最大积分曲线$\gamma_p(t)$可以定义
一个变换{\footnote{映射一般定义成两个不同集合之间的关系，比如$\phi:M\to N$；
        如果是同一集合内部的映射（$\phi:M\to M$），
        则赋予它一个特别的名字：\textbf{变换}。}}
\eq{\label{chdm:eqn_tmp-gammap}
    \phi(t,p):=\gamma_p(t), \quad \forall t\in \R,\quad \forall p\in M .
}
当$p$固定时，$\phi(t,p)$就是过$p$点积分曲线。
当参数$t$固定时，$\phi(t,p)$给出一个从$M$到$M$的映射（变换）；
$\forall(t,p )\in\R\times M$，记$\phi_t(p):=\phi(t,p)$。
可证明此变换满足如下定理。

\begin{theorem}\label{chdm:thm_1PTG}
    变换$\phi$满足如下条件：
    (1) $\phi_0= \id : M\to M$；
    (2) 对于任意实数$s,t$都有$\phi_t\circ \phi_s = \phi_{t+s}$ 。
\end{theorem}
\begin{proof}
    (1)： $\phi_0(p)=\phi(0,p)$自然将$p$点映射成$p$点，所以它是恒等映射。

    (2)：由命题\ref{chdm:thm_st}可得：
    $\phi_t\circ \phi_s(p)=\phi_t\bigl(\phi_s(p)\bigr) =\phi(t,q)= \phi_{t+s}(p)$ 。
\end{proof}
满足定理\ref{chdm:thm_1PTG}中两个条件的$\phi$（由式\eqref{chdm:eqn_tmp-gammap}定义）
被称为\textbf{（左）单参数可微变换群}。
完备切矢量场$v$一般也称为由单参数可微变换群$\phi(t,p)$诱导的切矢量场；两者一一对应。

还需要说明$\phi_t$是一个群。把第(2)条的复合映射看成群乘法，很明显此乘法是封闭的。
由第(2)条易证群乘法满足结合律。
$\phi_0$是恒等映射，可看成单位元。
再从第(2)条可知$\phi_{-t}\circ \phi_t=\phi_0=\id$，
因此$\phi_{-t} = \phi_t^{-1}$；这说明$\phi_t$的逆存在且是$C^\infty$的。
综上，$\phi_t$是群。

%对于任意给定的$t\in \R$，用反证法可证明$\phi_t(p)$是单射。
%假设$\phi_t(p)$不是单射，即存在两个不同的点$p$和$q$使得$\phi(t,p)=\phi(t,q)$；
%依据条件(2)有$\phi(-t,p)\circ\phi(t,p)=\phi(-t,p)\circ\phi(t,q)
%{\Rightarrow} \id=\phi(-t,p)\circ\phi(t,q)$，因
%$p$与$q$是两个不同的点，此式右端不是恒等映射，矛盾，所以必然单射。
%既然是单射，而且其逆存在，那必然也是满的；自然是双射。

%那么，$\gamma_p(t): \R \to M$是流形$M$中经过点$p$的一条可微曲线，称为单参数可微
%变换群$\phi(t,p)$经过点$p$的\textbf{轨道}（orbit）。


因$\phi_t^{-1}$存在，所以$\phi_t$必然是双射，进而$\phi_t, \ \phi_t^{-1}$都是$C^\infty$的双射；
故可知它是\textbf{微分同胚}的。
由第(2)条可得：$\phi_s\circ \phi_t = \phi_{s+t}= \phi_{t+s}=\phi_t\circ \phi_s$，
这说明此群是可对易的。


$v_p$是积分曲线$\gamma_p(t)$在$p$的切矢量，由式\eqref{chdm:eqn_DiffCurve}有公式，有
\eq{\label{chdm:eqn_tmp-vecp}
    v_p(f)= \left. \frac{\d \bigl(f \circ \gamma_p (t)\bigr)}{\d t} \right |_{t=0}
    = \lim_{t\to 0} \frac{f\bigl(\phi(t,p)\bigr) - f\bigl(\phi(0,p)\bigr)}{t},
}
其中$\forall f\in C^\infty _p(M) $。
利用上式，用推前映射语言来描述一下由$\phi(t,p)$诱导的切矢量场$v$的变换关系。
    很明显$p=\phi(0,p)$，令$q:=\phi(s,p)$。
    由式\eqref{chdm:eqn_tmp-vecp}易得
    \begin{align*}
        v_q(f)=& \lim_{t\to 0} \frac{f\bigl(\phi(t+s,p)\bigr) - f\bigl(\phi(s,p)\bigr)}{t}
         =\lim_{t\to 0} \frac{f\bigl(\phi_s(t,p)\bigr) - f\bigl(\phi_s(0,p)\bigr)}{t} \\
         =&\lim_{t\to 0} \frac{f\circ \phi_s(t,p) - f\circ \phi_s(0,p)}{t}
         =v_p\bigl(f\circ \phi_s(p) \bigr)  =\phi_{s*}(v_p)(f).
    \end{align*}
    因$f$及$p$的任意性，由上式可得
    \eq{\label{chdm:eqn_phisv}
        v_q = v_{\phi_s(p)} = \phi_{s*}(v_p) .
    }
    这便是两点间切矢量推前映射关系式。
%\end{example}


%设流形$M$有局部坐标系$(U;x^i)$，则由$p$映射出的像点$\phi(t,p)$的坐标表达式为
%\eq{\label{chdm:eqn_tmp-gxerp}
%    x^i=\phi^i(t,x^1_p,\cdots,x^m_p),
%}
%根据定义，$\phi^i$是所有参量$t,x^1_p,\cdots,x^m_p$的$C^\infty$函数。
%借用式\eqref{chdm:eqn_v-expand-on-natural-bases}，将切矢量场$v$在局部坐标系展开，
%\eq{\label{chdm:eqn_tmp-va2dfd}
%    v^i \left(\frac{\partial }{\partial x^i}\right)_{p}  \equiv v \
%    = \left. \frac{\partial \phi^i(t,x^1_p,\cdots,x^m_p)}{\partial t} \right|_{t=0}
%    \left(\frac{\partial }{\partial x^i}\right)_{p}
%    = \left. \frac{\partial x^i}{\partial t} \right|_{t=0}
%    \left(\frac{\partial }{\partial x^i}\right)_{p}
%}
%很容易看到它的系数是$C^\infty$的，所以这个矢量场是光滑的。
%由式\eqref{chdm:eqn_tmp-va2dfd}可以看到矢量场$v$在局部坐标系的
%系数$v^i=\frac{\partial \phi^i(t,x^1_p,\cdots,x^m_p)}{\partial t}|_{t=0}$是直接由
%映射$\phi$的对应分量求导得到；我们还可以把式\eqref{chdm:eqn_tmp-va2dfd}看成
%一常微分方程组，它有$m$个未知变量$x^i$，自变量是$t$，
%方程组可以简化为（仅在局部坐标上）
%
%式\eqref{chdm:eqn_tmp-va2dfd}非常清楚地表明了
%诱导切矢量场$v$是由轨道\eqref{chdm:eqn_tmp-gammap}对$t$求导得到的；
%那么轨道\eqref{chdm:eqn_tmp-gammap}自然是方程\eqref{chdm:eqn_tmp-va2xit}的解，
%也就是诱导切矢量场的积分曲线。
%
%
%我们在上面的论述中说明了轨道会诱导出流形$M$的切矢量场；反之，如果
%给定流形$M$上的一个切矢量场$v$是否存在$M$上的单参数可微变换群$\phi_t$，
%它以$v$为其诱导切矢量场？回答是肯定的！
%无非是解方程\eqref{chdm:eqn_tmp-va2xit}得到其积分曲线，再证明积分曲线
%局部单参数可微变换群；不过论证过程略显麻烦，为此先引入如下定义。

上面定义是针对整个流形$M$而言，它还可以弱化，只针对局部子集$U$，
\begin{definition}\label{chdm:def_1PTG-local}
    设$U$是光滑流形$M$的一个开子集，取$\epsilon>0$。
    如果光滑映射$\phi:(-\epsilon,\epsilon)\times U \to M $满足如下条件：

    (1) $\forall p\in U,\ |t|< \epsilon$都有$\phi_0 : U\to U$恒等映射；

    (2) 若$|s|<\epsilon,\ |t|<\epsilon, \ |s+t|<\epsilon$，则$\phi_s\circ \phi_t = \phi_{s+t}$ 。

    那么，称$\phi_t$为作用在开子集$U$上的\textbf{局部（左）单参数可微变换群}。
\end{definition}
局部单参数可微变换群自然也有与之对应的局部光滑诱导切矢量场$v\in \mathfrak{X}(U)$。
所诱导的切矢量场$v$一般说来不再是完备的，但$\phi_t$仍是局部微分同胚的。
需要强调的是：本节定义的单参数可微变换群是{左作用}在流形$M$上的；
与之对偶，可以定义{右作用}在流形$M$上的单参数可微变换群；
具体参见\S\ref{chlg:sec_Lie-transformation-group}。


%
%局部单参数可微变换群$\phi_t$同样可以诱导开子集$U$上的一个光滑切矢量场。$\forall p\in U$，
%取此点的局部坐标系为$(V;x^i)$，并且$V$是$U$的开子集。由于$\phi_t$是$C^\infty$的，所以
%肯定是连续映射，只要$\epsilon$充分小，总有$\phi_t(p)\in V$，于是$\phi_t$的诱导切矢量场
%的局部坐标可以表示为
%\eq{
%    v_p= \sum_{i=1}^{m} v^i_p \left(\frac{\partial }{\partial x^i}\right)_{p},\quad
%    {\text{其中}}\ v^i_p=\left.\frac{\partial \phi^i(t,x^1_p,\cdots,x^m_p)}{\partial t}\right|_{t=0}.
%}
%仿照前面表述，上式用微分方程表示便是式\eqref{chdm:eqn_tmp-va2xit}（注意$x^i\equiv\phi^i$）。
%
%\begin{theorem}\label{chdm:thm_vphil}
%    设$v$是光滑流形$M$上的光滑切矢量场，则$\forall p\in M$存在一个开邻域$U\subset M$以及
%    作用在$U$上的局部单参数变换群$\phi_t$，当$|t|<\epsilon$时，使得$v|_U$是$\phi_t$在$U$上
%    诱导的光滑切矢量场。
%\end{theorem}
%\begin{theorem}\label{chdm:thm_vphig}
%    设$v$是紧致的$m$维光滑流形$M$上的光滑切矢量场，
%    那么$v$可以生成作用在流形$M$上的单参数可微变换群$\phi:\R\times M\to M$，
%    即$v$是$\phi$诱导的光滑切矢量场。
%\end{theorem}
%\begin{proof}
%    以上两个定理证明请参阅文献\parencite[\S 6.2]{cc2001-zh}或\parencite[\S 3.3]{chenwh2001}。
%    需要提醒的是：单参数可微变换群$\phi$就是切矢量场$v$的积分曲线。
%\end{proof}
%\begin{remark}\label{chdm:rmk_vphilg}
%    在满足李普希兹条件下（通常都满足），常微分方程组解存在且唯一；故上面两定理中
%    的单参数可微变换群$\phi$是唯一确定的。定理\ref{chdm:thm_vphil}仅在局部成立；
%    而定理\ref{chdm:thm_vphig}在整个流形$M$上成立。%，其中条件“$M$是紧致”不可缺少。
%\end{remark}


\begin{theorem}\label{chdm:thm_1PDG-diff}
    设$\phi_t$是作用在光滑流形$M$上的单参数可微变换群，$v$是其所诱导的
    光滑切矢量场。如果$\psi:M\to M$是一个$C^\infty$微分同胚，
    那么$\tilde{\phi}_t := \psi \circ \phi_t \circ \psi^{-1}$也是作用在$M$上的单参数可微变换群，
    并且其诱导的光滑切矢量场是$\psi_{*}v$。
\end{theorem}
\begin{proof}
    显然$\tilde{\phi}_t$是从$M$到$M$的变换；$\tilde{\phi}_0=\psi \circ \phi_0 \circ \psi^{-1}=\id$是
    恒等映射，满足定理\ref{chdm:thm_1PTG}中的第(1)条；
    $\tilde{\phi}_t\circ \tilde{\phi}_s =  (\psi \circ \phi_t \circ \psi^{-1} ) \circ
    (\psi \circ \phi_s \circ \psi^{-1}) = \psi \circ \phi_{t+s} \circ \psi^{-1}=\tilde{\phi}_{t+s}$，
    也满足第(2)条；综上，$\tilde{\phi}_t$是作用在$M$上的单参数可微变换群。

    对其求导便可得诱导切矢量场，令$q=\psi(p)$，$\forall f\in C^\infty (M)$，有
    \begin{align*}
        \tilde{v}(f) & = \left. \frac{\d \bigl(f \circ \tilde{\phi}(t,q)\bigr)}{\d t} \right |_{t=0}
        =\left. \frac{\d \bigl(f \circ \psi \circ \phi_t \circ \psi^{-1} (q)\bigr)}{\d t} \right |_{t=0}
        =\left. \frac{\d \bigl(f \circ \psi \circ \phi(t,p))\bigr)}{\d t} \right |_{t=0} \\
        &\xlongequal{\ref{chdm:eqn_tmp-vecp}} v(f \circ \psi)
        \xlongequal{\ref{chdm:eqn_pushforward-vector-onfield}} \psi_{*}v(f) .
    \end{align*}
    因此$\tilde{\phi}_t$诱导的光滑切矢量场是$\tilde{v}=\psi_{*}v$。
\end{proof}

\begin{definition}\label{chdm:def_invariant-vector-field}
    设$v$是光滑流形$M$上的光滑切矢量场，$\psi:M\to M$是光滑微分同胚。
    如果$\psi_{*}v=v$，那么称$v$在$\psi$作用下是不变的。
\end{definition}


\begin{proposition}\label{chdm:thm_ivf}
    设$\phi_t$是作用在光滑流形$M$上的单参数可微变换群，$v$是其所诱导的
    光滑切矢量场；$\psi:M\to M$是一个$C^\infty$微分同胚。
    $v$在$\psi$作用下是不变的充分必要条件是：$\phi_t$与$\psi$可交换，
    即$\phi_t\circ \psi = \psi \circ \phi_t$。
\end{proposition}
\begin{proof}
    先证“$\Leftarrow$”。
    如果$\phi_t$与$\psi$可交换，即${\phi}_t = \psi \circ \phi_t \circ \psi^{-1}$，
    由定理\ref{chdm:thm_1PDG-diff}可知${v}=\psi_{*}v$，即$v$在$\psi$作用下是不变的。

    再证“$\Rightarrow$”。
    由${v}$和$\psi_{*}v$决定的
    单参数可微变换群分别是${\phi}_t$和$\psi \circ \phi_t \circ \psi^{-1}$。
    如果$v$在$\psi$作用下是不变的，即${v}=\psi_{*}v$，那么说明两个单参数可微变换群
    满足同一常微分方程和相同的初始条件，也就是两条积分曲线是重合的，
    即${\phi}_t = \psi \circ \phi_t \circ \psi^{-1}$。
\end{proof}


\begin{theorem}\label{chdm:thm_xy-yx}
    设$\phi_t$是作用在光滑流形$M$上的单参数可微变换群，$X$是其所诱导的光滑切矢量场；
    $Y$是$M$上另一切矢量场。那么两个切矢量场的Poisson括号（对易子）表达式为
    \eq{\label{chdm:eqn_xy-yx}
        \left[X,Y\right] = \lim_{t\to 0} \frac{Y_{\phi_{0}} - (\phi_{t})_{*} Y_{\phi_{-t}} }{t}
        =\lim_{t\to 0} \frac{(\phi_{-t})_{*} Y_{\phi_{t}} - Y_{\phi_{0}}}{t} .
    }
\end{theorem}
\begin{proof}
    只需令$t\to -t$即可看到两个表达式等价，所以只需证明其中之一便可。

    对于任意的$f\in C^\infty(M)$，有$(\phi_{-t})_{*} Y_{\phi_{t}}(f)=Y_{\phi_{t}}\bigl( f\circ\phi({-t},p) \bigr)$，
    其中点$p$（对应$t=0$点）是$M$中任一点。从中可以看出，需要仔细计算其中的
    表达式$F(t):= f\circ\phi({-t},p)$；在$p\in M$点小邻域内，我们利用如下积分技巧
    \begin{equation*}
        F(t)-F(0) = \int_{0}^{1} \frac{\d F(st)}{\d s} \d s
        = t \cdot \int_{0}^{1}  \frac{\d F(st)}{\d (st)} \d s
        =t\cdot g(t,p).
    \end{equation*}
    记上式积分为 $g(t,p)$，很明显$g(t,p)$是光滑函数。由此得
    \eq{
        f\circ\phi({-t},p) = f(p) + t \cdot g(t,p) .
    }
    先计算一下$g(0,p)$
    \begin{align*}
        g(0,p)=& \lim_{t\to 0} \int_{0}^{1}  \frac{\d F(st)}{\d (st)} \d s
        = \int_{0}^{1} \lim_{t\to 0} \frac{\d f\circ\phi({-st},p)}{\d (st)} \d s \\
        \xlongequal{\ref{chdm:eqn_tmp-vecp}} &
        - X_{p}(f) \int_{0}^{1}  \d s = -X_{p}(f) .
    \end{align*}
    再计算极限（注意应用上式）
    \begin{align*}
        &\lim_{t\to 0} \frac{(\phi_{-t})_{*} Y_{\phi_{t}} - Y_{\phi_{0}}}{t}(f) =
        \lim_{t\to 0} \frac{Y_{\phi_{t}}\bigl( f(p) + t \cdot g(t,p) \bigr) - Y_{\phi_{0}}(f)}{t}  \\
        &= \lim_{t\to 0} \frac{Y_{\phi_{t}}\bigl( f(p)  \bigr) - Y_{\phi_{0}}(f)}{t}
          + \lim_{t\to 0} Y_{\phi_{t}}\bigl(  g(t,p) \bigr) \\
        &=\lim_{t\to 0} \frac{(Y_{p} f) \circ\phi(t,p)  - (Y_{p}f )\circ\phi(0,p) }{t} - Y_{p} \circ X_{p}(f) \\
%        &=\left. \frac{\d Y_{p}\bigl( f \circ\phi(t,p) \bigr)}{\d t}\right|_{t=0} - Y_{p} \circ X_{p}(f) \\
        & \xlongequal{\ref{chdm:eqn_tmp-vecp}}
        X_{p}\Bigl(Y_{p}\bigl( f (p) \bigr)\Bigr) - Y_{p} \circ X_{p}(f)
        =\left[X,Y\right]_p (f) .
    \end{align*}
    因$p$是任选的，所以上式对整个流形均正确。
\end{proof}


\begin{theorem}\label{chdm:thm_1PDG-ppy}
    设$X$是光滑流形$M$上的光滑切矢量场，如果在点$p \in M$矢量场$X_p \neq 0$，
    那么，存在$p$点的一个局部坐标系$(V;y^i)$使得$X |_V = (\frac{\partial }{\partial y^1} )$。
\end{theorem}
\begin{proof}
    设$p$点附近存在局部坐标系$(U;x^i)$，并将$p$取为坐标原点，即$x^i(p)=0$。在这个局部坐标系内
    光滑流形$X$可以展开为：$ X |_U = \xi^ j (\frac{\partial }{\partial x^j} )$，
    其中 $\xi^ j \in C^\infty(U)$。
    因$X_p \neq 0$，不妨设$\xi^1(p)\neq 0$。利用$\xi^j$的连续性，可假设$\xi^1$在
    开集$U$上不为零；如果有零点，将$U$缩小便可。
    考虑切矢量场$X$的积分轨道，要求此轨道经过坐标面$x^1=0$上的点$(0,y^2,\cdots,y^m)$，
    此轨道满足如下常微分方程和初始条件：
    \begin{equation*}
        \frac{\d x^i}{\d t} = \xi^i (x^1,\cdots,x^m); \quad
        x^1(0;y^2,\cdots,y^m) = 0; \quad
        x^\alpha(0;y^2,\cdots,y^m) = y^\alpha.
    \end{equation*}
    其中$2 \leqslant \alpha \leqslant m$。
    根据常微分方程理论（可参阅文献\parencite[\S 31]{arnold-2001-ode}），上述方程
    的解$x^i(t;y^2,\cdots,y^m)$光滑依赖于自变量$t$和初始值$y^2,\cdots,y^m$。
    为了行文方便，令$y^1=t$，解的Jacobi行列式为
    \begin{equation*}
        \left. \frac{\partial (x^1,\cdots,x^m)}{\partial (y^1,\cdots,y^m)} \right|_{y=0}=
        \begin{vmatrix}
            \xi^1(p) & 0 &\cdots & 0 \\
            \xi^2(p) & 1 &\cdots & 0 \\
            \vdots   & \vdots &\ddots & \vdots \\
            \xi^m(p) & 0 &\cdots & 1 \\
        \end{vmatrix}
        = \xi^1(p) \neq 0.
    \end{equation*}
    Jacobi行列式不为零，说明可把坐标$(x^1,\cdots,x^m)$变换成另一套局部
    坐标$(y^1,\cdots,y^m)$，即存在$p$点邻域$V\subset U$的局部坐标系$\{y\}$。
    求复合函数偏导数，
    \begin{equation*}
        \frac{\partial}{\partial y^1} = \frac{\partial x^i}{\partial y^1} \frac{\partial}{\partial x^i}
        = \xi^i \frac{\partial}{\partial x^i}
    \end{equation*}
    这就是所求的坐标表达式。
    光滑矢量场$X$存在与之对应的单参数可微变换群$\phi(t,p)$；由上易见，
    单参数$t$便是所选坐标系的第一个坐标。
    一般称坐标系$(V;y^i)$是矢量场$X$的\textbf{适配坐标系}。
\end{proof}

\section{Lie 导数}\label{appx:Lie-deri}

进一步对微分同胚，则称\textbf{单参微分同胚群}。对 $\forall p\in M$，取不同参数 $t$ 的 $\phi_t$ 将其映射至不同点 $\phi_t(p)$，形成一条以 $\phi_0(p)=p$ 为参数零点的光滑曲线 $C_p:\R\to M$，这称为该单参微分同胚群过 $p$ 点的\textbf{轨道}(orbit)。

因此 $M$ 上的单参微分同胚群给出其对应的光滑切矢场。

由于同胚是双射，因此各轨道是不会相交的，因此只要给出各 $p$ 点在其对应轨道的切矢 $\pdv*{t}$，就给出了 $M$ 上的光滑切矢场。

反之是否成立？可以证明，


看来 $\forall \sigma\in\R$ 可借用矢量场 $V$ 的积分曲线定义微分同胚 $\phi_\sigma:M\to M$：设过 $p$ 的积分曲线为 $C_p$（且 $C_p(0)=p$），规定 $\phi_\sigma(p):=C_p(\sigma)$，于是似乎得到了单参微分同胚群。然而存在极端情形，某条积分曲线在参数取某些值时像点不存在，比如可人为挖去背景的某些区域，因此只能说 $V$ 局部地给出单参微分同胚群。准确地，只要 $p\in M$ 存在，总可找到其开邻域 $U\subset M$，使得 $U,\phi_t[U]$ 同胚且满足$\phi_t\phi_s=\phi_{t+s}$，其中参数范围为零的某开邻域 $I\subset\R$。这样我们就把 $\{\phi_t|t\in I\}$ 称为 $p$ 的\textbf{单参微分同胚局部群}。

一般来说总考虑 $U=M$。

数学上，如果一条曲线处处与该矢量场相切，就称为该矢量场的一条积分曲线。十分显然，因为一点只对应一条积分曲线，倘若对应多条，矢量场于该点的取值便出现多值。可见积分曲线不会相交。

\textbf{参考系}指时空上的一个光滑单位类时矢量场 $\bm Z$。
参考系的每一积分曲线 $\gamma(\tau)$ 都是类时世界线。

所有微分同胚（包括那些不同单参变换的复合，它们不是单参变换）仍将构成群，称为 $M$ 的\textbf{微分同胚群}。一个集合可以是群，也可同时是一个背景，这种群称为 \textbf{Lie 群}。连续变换构成的群就可以用参数代表其坐标系，故实质上 Lie 群，其维数即参数种类，而一个切矢场又给出一种参数，可见微分同胚群是一种无穷维 Lie 群。一个光滑切矢场生成的是微分同胚群的一个\textbf{单参子群}，故又被称为微分同胚群的一个\textbf{无穷小生成元}或\textbf{生成元}。

矢量场 $\xi$ 能局部地给出单参微分同胚群 $\{\phi_t\}$。由 $\phi_t$ 可诱导出一种导数，其目的是为了比较邻近的 $T,\phi_{t*} T$ 或等价的 $\phi_t^*T,T$。
\begin{definition}
    张量场沿切矢场 $\xi$ 的 \textbf{Lie 导数}为
\eq{
\Lie_\xi T:=\lim_{t\to 0}\frac{\phi_t^* T-T}{t}.
}
\end{definition}

\begin{example}
    就标量场 $f$ 而言 $\Lie_\xi f=\xi(f)$。
\end{example}


设$X$是$m$维光滑流形$M$上的完备光滑切矢量场，它诱导的（局部）单参数可微变换群是$\phi_t$，
经过点$x\in M$的轨道记为$\gamma_x(t)$（当$t=0$时$\gamma_x(0)=x$）；
由于$\phi_t$是（局部）微分同胚，可将轨道$\gamma_x(t)$上不同点的切空间建立联系。
\begin{definition}
    由 $\phi_t$ 可定义 Lie 导数\footnote{Sophus Lie（1842--1899），挪威数学家。其另一重大贡献是李群，见第\ref{chlg}章。}：
    \begin{align}
        (\Lie_Xf)(p)&:=\lim_{t\to0}\frac{f(\phi_t(p))-f(p)}{t}=X_p(f),
        \label{chdm:def_LieD-scalar}\\
        \Lie_XY&:=\lim_{t\to0}\frac{(\phi_{-t})_*Y-Y}{t},
        \label{chdm:def_LieD-vector}\\
        \Lie_XT&:=\lim_{t\to0}\frac{\phi_t^*T-T}{t}.
        \label{chdm:def_LieD-tensor}
    \end{align}
    其中 $f\in C^\infty(M)$、$X,Y\in\mathfrak X(M)$、$T\in\mathfrak T^p_q(M)$。这些极限均逐点在同一张量空间内求取，前两式是第三式的特例。

\end{definition}


由于Lie 导数定义并不十分明朗，下面略微详细描述一下。
由于$\phi_t$（至少局部上）是微分同胚，它的诱导映射都是线性同构的。
\begin{align}
    \phi_{-t*} = \phi_{t*} ^{-1} &: T_{\gamma_x(t)}M \to T_x M , \\
    \phi_{t}^{*} &: T_{\gamma_x(t)}^{*}M \to T_x^{*} M , \\
    \phi_{t}^{*} &: T_q^p\bigl({\gamma_x(t)}\bigr) \to T_q^p(x) .
\end{align}
上式最后一行的$\phi_t^{*}$是点$x$张量空间和点$\gamma_x(t)$张量空间的线性同构映射。
设有局部坐标系$(U;x^i)$，并令$t$是个充分小的正实数，则有$\gamma_x(t)\in U$；
我们用$\{x^i\}$表示$x$点处的坐标分量，为了在形式上有所区分，我们用$\{y^i\}$表示$\gamma_x(t)$点处的坐标分量；
需要注意的是，当$t\to 0$时，坐标$\{y^i\}$就变成了坐标$\{x^i\}$，
即$\lim_{t\to 0} y^i = x^i $。

记 $y=\gamma_x(t)$，并以 $e_i(x),e_i(y)$ 及 $\epsilon^j(x),\epsilon^j(y)$ 表示同一坐标系在两点的切基及对偶基。则
\[
    T_y=T^{i_1\cdots i_p}{}_{j_1\cdots j_q}(y)
    e_{i_1}(y)\otimes\cdots\otimes e_{i_p}(y)\otimes
    \epsilon^{j_1}(y)\otimes\cdots\otimes\epsilon^{j_q}(y).
\]
拉回后各因子均位于 $x$ 处：
\eq{\label{chdm:eqn_tmp501}
\begin{aligned}
    (\phi_t^*T)_x={}&T^{i_1\cdots i_p}{}_{j_1\cdots j_q}(y)
    \bigl((\phi_{-t})_*e_{i_1}(y)\bigr)\otimes\cdots\otimes
    \bigl((\phi_{-t})_*e_{i_p}(y)\bigr)\\
    &\otimes\phi_t^*\epsilon^{j_1}(y)\otimes\cdots\otimes\phi_t^*\epsilon^{j_q}(y).
\end{aligned}
}
因此 $\lim_{t\to0}(\phi_t^*T)_x=T_x$。上述表达式与坐标选择无关：坐标变化时，分量与基底的变换矩阵互相抵消。

\begin{proposition}\label{chdm:thm_cp-pc}
    拉回 $\phi_t^*$ 与缩并运算 $C$ 可交换。
\end{proposition}
\begin{proof}
    记 $y=\phi_t(x)$。对 $v\in T_yM$、$\omega\in T_y^*M$，由拉回定义，
    \[
        (\phi_t^*\omega)\bigl((\phi_{-t})_*v\bigr)
        =\omega\bigl((\phi_t)_*(\phi_{-t})_*v\bigr)=\omega(v).
    \]
    缩并即对一个切矢因子和一个余切矢因子作此配对，故对简单张量有
    \eq{\label{chdm:eqn_tmpa852}
        C(\phi_t^*T)=\phi_t^*(CT).
    }
    任意张量均为简单张量之和，由线性性即得结论。在坐标基下，上述配对保持性也可由互逆的 Jacobi 矩阵看出：
    \eq{\label{chdm:eqn_tmp-transxy}
        (\phi_{-t})_*\p_{y^\alpha}=\pdv{x^i}{y^\alpha}\p_{x^i},\quad
        \phi_t^*(\d y^\alpha)=\pdv{y^\alpha}{x^i}\d x^i.
    }
\end{proof}

\begin{theorem}\label{chdm:thm_LieD-property}
    Lie 导数$\Lie_{X}$有如下运算规则（其中，$\forall T\in \mathfrak{T}^p_q(M),\ S\in \mathfrak{T}^r_s(M)$，
    $\forall \lambda \in \R$， $\forall \omega\in \mathfrak{X}^*(M)$，
    $\forall X,Y\in \mathfrak{X}(M)$；
    其中加法要求 $T,S$ 同型）：

    (1) 线性：$\Lie_{X}(T + \lambda \cdot S) = \Lie_{X}T + \lambda\cdot \Lie_{X}S$；

    (2) Leibniz律：$\Lie_{X} (T \otimes S) = T \otimes \Lie_{X} S + (\Lie_{X} T )\otimes S$；

    (3) 与缩并运算可交换，即$C (\Lie_{X} T) =  \Lie_{X} (C T) $；


    (4) 若$\psi:M\to M$是微分同胚，则$\psi_*(\Lie_{X} T) = \Lie_{\psi_*X} \psi_* T$。
\end{theorem}
\begin{proof}
    (1) 根据定义式\eqref{chdm:def_LieD-tensor}很容易证明线性性。

    (2) 根据定义式\eqref{chdm:def_LieD-tensor}和命题\ref{chdm:thm_pptt}，有
    \begin{align*}
        \Lie_{X} (T \otimes S) & = \lim_{t\to 0} \frac{\phi^{*}_{t} (T\otimes S) -T\otimes S }{t}
          = \lim_{t\to 0} \frac{\phi^{*}_{t} (T)\otimes \phi^{*}_{t}(S) -T\otimes S }{t}\\
        & = \lim_{t\to 0} \frac{\phi^{*}_{t} (T)\otimes \phi^{*}_{t}(S)
             -T\otimes \phi^{*}_{t}(S) +T\otimes \phi^{*}_{t}(S) -T\otimes S }{t} \\
        & = T \otimes \Lie_{X} S + (\Lie_{X} T )\otimes S .
    \end{align*}

    (3) 根据定义式\eqref{chdm:def_LieD-tensor}和命题\ref{chdm:thm_cp-pc}非常容易证明此命题；留给读者当练习。


    (4) 依据定理\ref{chdm:thm_1PDG-diff}可令$\tilde{\phi}_t := \psi \circ \phi_t \circ \psi^{-1}$；
    先验证切矢量场公式
    \begin{equation*}
        \psi_*(\Lie_{X} Y)  = \lim_{t\to 0} \frac{\psi_*{\phi}_{-t*} Y - \psi_*Y }{t}
       = \lim_{t\to 0} \frac{\tilde{\phi}_{-t*} \psi_*Y -\psi_*Y }{t}
       =\Lie_{\psi_*X} \psi_*Y .
    \end{equation*}
    再验证余切矢量场公式
    \begin{align*}
        \psi^{-1*}(\Lie_{X} \omega)  = \lim_{t\to 0} \frac{\psi^{-1*} {\phi}_{t}^{*}
            \omega - \psi^{-1*}\omega }{t}
        = \lim_{t\to 0} \frac{ \tilde{\phi}_{t}^{*} \psi^{-1*}\omega - \psi^{-1*}\omega }{t}
        =\Lie_{\psi_*X} \psi^{-1*}\omega  .
    \end{align*}
    利用Leibniz律，可推广至高阶张量场。
\end{proof}

\begin{example}
    切矢量场、余切矢量场的 Lie 导数
\end{example}
由定理\ref{chdm:thm_xy-yx}可知切矢量场的 Lie 导数可以表示为：
\eq{\label{chdm:eqn_LieD-tangent}
    \Lie_{X} Y = [X, Y] , \quad \forall X, Y \in \mathfrak{X}(M).
}
利用式\eqref{chdm:def_LieD-scalar}及 Lie 导数与缩并可交换，得
\[
    (\Lie_X\omega)(Y)+\omega(\Lie_XY)=\Lie_X\bigl(\omega(Y)\bigr)=X\bigl(\omega(Y)\bigr).
\]
于是余切矢量场的 Lie 导数满足
\eq{\label{chdm:eqn_LieD-cotangent}
    (\Lie_X\omega)(Y)=X\bigl(\omega(Y)\bigr)-\omega([X,Y]),\quad X,Y\in\mathfrak X(M).
}
一般张量场可通过作用于多个切矢场和余切矢场来描述。在引入无挠联络后，还有分量表达式，见\S\ref{chccr:sec_LieD-Nabla}。


\begin{example}
    试证明：$\Lie_{X+Y} T = \Lie_{X} T  +\Lie_{Y} T $；其中$T$是任意张量场。
\end{example}
由式\eqref{chdm:eqn_LieD-tangent}、\eqref{chdm:eqn_LieD-cotangent}出发，
不难验证本例题对切矢量场、余切矢量场是正确的。再利用Leibniz律可将本例推广到高阶张量场。
需要注意，一般说来$\Lie_{f\cdot X} T \neq f\cdot \Lie_{X} T$，其中$f\in C^\infty(M)$；
此时需用式\eqref{chccr:eqn_LieD-tensor-Nabla}来计算。


\begin{example}
    适配坐标系中的 Lie 导数。
\end{example}

由定理\ref{chdm:thm_1PDG-ppy}，若 $X_p\ne0$，可取适配坐标系使 $X|_U=\p_1$。此时流只改变第一个坐标：
\[
    y^1=x^1+t,\quad y^k=x^k\quad(k>1),
\]
故其 Jacobi 矩阵为单位矩阵，坐标基及对偶基在流的推前和拉回下不变。以 $(1,1)$ 型张量场 $T$ 为例，由式\eqref{chdm:eqn_tmp501}，
\[
    (\phi_t^*T)_p=T^i{}_j(\phi_t(p))\p_i|_p\otimes\left.\d x^j\right|_p.
\]
因此
\eq{\label{chdm:eqn_LieD-adaptedX}
    (\Lie_XT)^i{}_j(p)
    =\lim_{t\to0}\frac{T^i{}_j(\phi_t(p))-T^i{}_j(p)}{t}
    =\p_1T^i{}_j(p)=X(T^i{}_j)(p).
}
这里 $(\Lie_XT)^i{}_j=(\Lie_XT)(\d x^i,\p_j)$ 是张量场 $\Lie_XT$ 的分量，而 $X(T^i{}_j)$ 是对分量函数求导；二者在适配坐标系下相等，在一般坐标系下未必相等。


\begin{exercise}
	补全证明式\eqref{chdm:eqn_LieD-adaptedX}中所省略的步骤。
\end{exercise}

\begin{exercise}
	在二维流形$\R^2$上有局部坐标系$x,y$，令$X = x^2 (\frac{\partial}{\partial x})$、
	$Y = (\frac{\partial}{\partial x}) + y^2 (\frac{\partial}{\partial y})$。
	计算$\Lie_X Y$关于$(\frac{\partial}{\partial x})$、$(\frac{\partial}{\partial y})$的表达式。
	作变换$u=x+y$、$v=x-y$，求$\Lie_X Y$关于$(\frac{\partial}{\partial u})$、$(\frac{\partial}{\partial v})$的表达式。
\end{exercise}
